%==============================================================================
% NeurIPS 2026 主控文件 — EntroDiff
%
% 提交规范严格遵守 paper/black/CONVENTIONS.md 与项目 CLAUDE.md：
%   - 9 页正文上限（含图表）
%   - 双盲：禁加 [final] / [preprint]
%   - 公式必须用 equation / align 环境，禁用 $$
%   - 表格用 booktabs，禁用垂直分割线
%   - 图表 caption 必须含 "Key take-away message"
%   - 引用 natbib，文内统一 \citet / \citep
%
% 编译：pdflatex → bibtex → pdflatex → pdflatex
%==============================================================================
\documentclass{article}

% natbib 选项必须在加载 neurips_2026 之前传入
\PassOptionsToPackage{numbers, compress}{natbib}

% 双盲提交：使用 default（= main + 双盲 + 行号）。严禁加 [final] 或 [preprint]
\usepackage{neurips_2026}

%------------------------------------------------------------------------------
% 通用宏包（与官方模板对齐）
%------------------------------------------------------------------------------
\usepackage[utf8]{inputenc}
\usepackage[T1]{fontenc}
\usepackage{hyperref}
\usepackage{url}
\usepackage{booktabs}    % 高质量表格（禁用 vertical rules）
\usepackage{amsfonts}
\usepackage{amsmath}     % equation / align / 等环境
\usepackage{amssymb}
\usepackage{amsthm}      % 定理 / 命题 / 引理环境
\usepackage{mathtools}
\usepackage{nicefrac}
\usepackage{microtype}
\usepackage{xcolor}
\usepackage{graphicx}
\graphicspath{{figures/}}

%------------------------------------------------------------------------------
% 项目自定义宏（专业符号统一）
%------------------------------------------------------------------------------
%==============================================================================
% macros/notation.tex — 统一项目符号表
%
% 设计原则：
%   1. 与 Docs/black/path_A_method_skeleton.md §3.1 Notation 表保持一致
%   2. 加宏的目的：(a) 一处改名全文同步 (b) 防止 reviewer 抱怨符号漂移
%   3. 所有专业符号都通过宏，不直接写裸 LaTeX
%
% 修改时务必同步更新 path_A_method_skeleton.md
%==============================================================================

% ----- 数集 -----
\newcommand{\R}{\mathbb{R}}
\newcommand{\N}{\mathbb{N}}
\newcommand{\Z}{\mathbb{Z}}
\newcommand{\E}{\mathbb{E}}
\newcommand{\Prob}{\mathbb{P}}

% ----- 概率 / 测度 -----
\newcommand{\Law}{\mathrm{Law}}
\newcommand{\KL}{\mathrm{KL}}
\newcommand{\Wass}[1]{W_{#1}}                 % e.g. \Wass{1}, \Wass{2}
\newcommand{\TV}{\mathrm{TV}}
\newcommand{\BV}{\mathrm{BV}}
\newcommand{\Ent}{\mathrm{Ent}}
\newcommand{\Lip}{\mathrm{Lip}}

% ----- 分布与路径 -----
\newcommand{\rhotrue}{\rho^{\star}}             % 目标 PDE 解分布
\newcommand{\rhotau}{\rho_{\tau}}               % 扩散时间 tau 处的分布
\newcommand{\ptau}{p_{\tau}}                    % noised density
\newcommand{\score}{s}                          % score
\newcommand{\sth}{s_{\theta}}                   % score network
\newcommand{\Dth}{D_{\theta}}                   % denoiser network
\newcommand{\sigtau}{\sigma(\tau)}              % 噪声尺度
\newcommand{\Gtau}{G_{\tau}}                    % heat kernel

% ----- PDE -----
\newcommand{\dt}{\partial_{t}}
\newcommand{\dx}{\partial_{x}}
\newcommand{\dtau}{\partial_{\tau}}
\newcommand{\Lap}{\Delta}
\newcommand{\flux}{f}                           % 通量函数 f(u)
\newcommand{\physvis}{\nu_{\mathrm{phys}}}      % PDE 物理黏性
\newcommand{\snr}{\mathrm{SNR}}
\newcommand{\amp}{\Lambda}                      % Score Shocks 的指数放大因子

% ----- 损失 -----
\newcommand{\Ldsm}{\mathcal{L}_{\mathrm{DSM}}}
\newcommand{\Lpde}{\mathcal{L}_{\mathrm{PDE}}}
\newcommand{\Lent}{\mathcal{L}_{\mathrm{ent}}}
\newcommand{\Lbv}{\mathcal{L}_{\mathrm{BV}}}
\newcommand{\Lburg}{\mathcal{L}_{\mathrm{Burg}}}
\newcommand{\Rent}{\mathrm{R}_{\mathrm{ent}}}

% ----- 集合 / 子集 -----
\newcommand{\Shockphys}{\Sigma_{\mathrm{phys}}}
\newcommand{\Shockscore}{\Sigma_{\mathrm{score}}}

% ----- 物理空间解与轨迹（SYMBOL.md §5/§10） -----
\newcommand{\physsol}{\mathbf{u}^{\star}}        % Kruzhkov 熵解
\newcommand{\physsolvis}{\mathbf{u}^{\nu}}       % 粘性 PDE 解
\newcommand{\initdat}{\mathbf{u}_{0}}            % PDE 初值
\newcommand{\rhotphys}{\rho_{t}}                 % 物理时间 t 处分布

% ----- Score 残差与误差 -----
\newcommand{\errsc}{e}                           % score 残差 e = s_theta - s

% ----- 几何对象扩展 -----
\newcommand{\dShock}{d_{\Sigma}}                 % 到 shock 流形的符号距离

% ----- BV-aware 参数化分量（SYMBOL.md §9） -----
\newcommand{\phisbg}{\phi^{\mathrm{sm}}_{\theta}}  % 平滑背景势函数
\newcommand{\phissh}{\phi^{\mathrm{sh}}_{\theta}}  % shock 符号距离场
\newcommand{\jumpamp}{\kappa_{\theta}}           % 局部跳跃强度

% ----- 轨迹与生成样本（SYMBOL.md §10） -----
\newcommand{\unat}{u^{\natural}}                 % 真实 score 驱动的反向 ODE 轨迹
\newcommand{\uhat}{\widehat{u}}                  % 网络 score 驱动的反向 ODE 轨迹
\newcommand{\uth}{u^{\theta}}                    % 最终生成样本 = \uhat(0)
\newcommand{\muth}{\mu_{\theta}}                 % 生成分布 = Law(u^theta)

% ----- 实用工具 -----
\newcommand{\abs}[1]{\lvert #1 \rvert}
\newcommand{\norm}[1]{\lVert #1 \rVert}

% ----- 排版小工具 -----
\newcommand{\eg}{e.g.,\ }
\newcommand{\ie}{i.e.,\ }
\newcommand{\cf}{cf.\ }

% ----- 备注高亮（写作期临时使用，提交前删除）-----
\newcommand{\todo}[1]{\textcolor{red}{\textbf{TODO: }#1}}
\newcommand{\note}[1]{\textcolor{blue}{[Note: #1]}}


%------------------------------------------------------------------------------
% 定理 / 命题环境
%------------------------------------------------------------------------------
\newtheorem{theorem}{Theorem}
\newtheorem{proposition}[theorem]{Proposition}
\newtheorem{lemma}[theorem]{Lemma}
\newtheorem{corollary}[theorem]{Corollary}
\theoremstyle{definition}
\newtheorem{definition}[theorem]{Definition}
\theoremstyle{remark}
\newtheorem{remark}[theorem]{Remark}

%------------------------------------------------------------------------------
% 标题与作者（双盲提交占位）
%------------------------------------------------------------------------------
\title{Shock-Aware Diffusion Solvers for Hyperbolic PDEs:\\A Double-Burgers Perspective}

% 双盲：占位作者；camera-ready 时再换上真实作者块
\author{%
  Anonymous Authors\\
  Anonymous Affiliation\\
  \texttt{anonymous@anonymous}\\
}

\begin{document}
\maketitle

\begin{abstract}
% TODO[abstract] · NeurIPS-style 4-sentence abstract:
%   1) Gap: state-of-the-art diffusion-based PDE solvers (DiffusionPDE, FunDPS) saturate on smooth PDEs but fail on hyperbolic / shock-containing equations.
%   2) Insight: the score field of a diffusion model itself satisfies a viscous Burgers equation; when the target PDE is also Burgers-type, two coupled Burgers structures emerge.
%   3) Method: we exploit this coupling via (i) viscosity-matched noise schedule, (ii) BV-aware score parameterization with an embedded tanh interfacial layer, (iii) entropy / BV / Burgers-consistency regularization.
%   4) Result: $W_1$ convergence rate $\mathcal O(\varepsilon^{1/2})$ to the Kruzhkov entropy solution (free of the $\exp(\Lambda)$ blow-up factor), validated on Burgers, Buckley–Leverett, and 1D Euler/Sod toy problems.
[Placeholder] Abstract to be drafted in W10. The four-sentence skeleton lives in this comment.
\end{abstract}

%------------------------------------------------------------------------------
%==============================================================================
% Section 1 · Introduction (~1 page)
%
% Narrative arc (NeurIPS 2026 标准三段式):
%   Gap:  diffusion PDE solvers fail at shocks — empirically, analytically, algorithmically
%   Insight:  Score Shocks reveals the score is a Burgers field → double-Burgers coupling
%   Contribution:  (C1) coupling theorem / (C2) tanh-hard-coded architecture / (C3) ε^{1/2} rate
%
% 写作原则: 每段结尾指向下一段 / 每句话有信息增量 / 不用"I/We/Our" (双盲匿名)
%==============================================================================
\section{Introduction}
\label{sec:intro}

% ----- Gap: the shock triple-blind-spot -----
Generative diffusion in function space is reshaping data-driven {PDE} solving. The current frontier---{DiffusionPDE} \citep{huang2024diffusionpde} and {FunDPS} \citep{yao2025fundps} in Banach spaces---achieves striking accuracy on Darcy, Helmholtz, and moderate-Reynolds Navier--Stokes. Yet hyperbolic conservation laws remain untouched. The avoidance is methodological, not accidental: this class is the simultaneous \emph{empirical}, \emph{analytical}, and \emph{algorithmic} blind spot of the paradigm.

\textit{Empirically}, the five {PDE} benchmarks of \citet{yao2025fundps} are every one elliptic or parabolic; the sole Burgers experiment of \citet{huang2024diffusionpde} fixes $\physvis = 10^{-2}$ to sidestep the inviscid limit. \textit{Analytically}, existing convergence theory for diffusion solvers operates in $L^{2}$, $\Wass{2}$, or $\TV$ under classical smoothness; none accommodates weak solutions, $L^{1}$-contraction, or the Kruzhkov entropy condition \citep{kruzhkov1970}. \textit{Algorithmically}, the {PDE}-residual guidance central to both {DiffusionPDE} and {FunDPS} evaluates spatial derivatives by central differences---a discretisation that is singular at discontinuity interfaces. The guidance velocity itself diverges at the very feature that defines the hyperbolic regime.

% ----- Insight: the double-Burgers pivot -----
A recent identity changes the calculus. \citet{sarkar2026scoreshocks} establish that the score field of a {VE-SDE} diffusion model satisfies a viscous Burgers equation along the diffusion-time axis: $\dtau \score + 2\score\cdot\nabla_{u}\score = \Lap_{u}\score$. Equivalently, $-2\score$ is itself a Burgers velocity field. When the target distribution consists of solutions to a hyperbolic conservation law---itself of Burgers type in space-time---two Burgers structures coexist in orthogonal directions, and they are \emph{coupled}: physical shocks define the singular support of the data law, thereby fixing the location of the score-level interfacial layer through the correspondence $\Shockphys(t) \leftrightarrow \Shockscore(\tau)$ (Figure~\ref{fig:double-burgers}, Theorem~\ref{thm:double-burgers}).

This coupling reframes shock geometry from an obstacle into an asset. In standard diffusion training, a neural denoiser must \emph{learn} the sharp $\tanh$ profile of the score near a shock from finite data; the approximation error there is amplified by an $\exp(\amp T_{d})$ factor \citep[Theorem~6.3]{sarkar2026scoreshocks}, rendering the shock region fragile. The double-Burgers picture suggests a different route: \emph{hard-code} the exact $\tanh$ interfacial profile into the network architecture, reducing the learning task from reproducing a singular function to merely locating it and estimating its amplitude.

% ----- Contributions -----
The following three contributions execute this programme and provide its theoretical certification (quantitative validation in \S\ref{sec:experiments}):

\begin{enumerate}[label=(\textbf{C\arabic*}), leftmargin=*, itemsep=2pt]
    \item \textbf{Double-Burgers coupling (Theorem~\ref{thm:double-burgers}).} We formalise the geometric correspondence between $\Shockphys(t)$ and $\Shockscore(\tau)$ as a coupled Burgers system on $(\tau, t)$, establishing the structural foundation.
    \item \textbf{BV-aware score parameterization (\S\ref{sec:method}).} The exact $\tanh$ interfacial profile of the score field \citep[Prop.~5.4]{sarkar2026scoreshocks} is embedded into the network as an architectural prior: $\sth = \nabla\phisbg + (\jumpamp/2)\,\tanh(\phissh/2\sigtau^{2})\,\nabla\phissh$. The network learns only the shock location ($\phissh$) and amplitude ($\jumpamp$), not the shock shape.
    \item \textbf{Improved rate $\Wass{1} \le \mathcal{O}(\varepsilon^{1/2})$ (Theorem~\ref{thm:improved-rate}).} Under (C2) and a $\BV$ regulariser, generated samples converge to the Kruzhkov entropy solution at rate $\varepsilon^{1/2}$ in the score-matching error, with constants \emph{independent} of the Score Shocks amplification exponent $\amp$.
\end{enumerate}

%==============================================================================
% Section 2 · Related Work and Background (~0.7 page)
%
% Structure: 4 related-work paragraphs (each ends with "where we differ") +
%            3 minimal background paragraphs (VE-SDE / hyperbolic / metrics)
%==============================================================================
\section{Related Work and Background}
\label{sec:related-bg}

We situate the paper in four research threads and provide minimal background for the hyperbolic and metric tools used throughout.

\paragraph{Function-space diffusion.}
\citet{lim2023ddo} pioneered score-based diffusion on Banach spaces. \citet{huang2024diffusionpde} introduced {PDE}-residual--guided sampling for forward and inverse problems. \citet{yao2025fundps} extended Tweedie sampling to infinite-dimensional Banach spaces, reaching $3\%$ observation density on five benchmarks---all of which are elliptic or parabolic. The common denominator is a {PDE}-residual term evaluated by central differences. \textbf{Where we differ:} we target \emph{hyperbolic} conservation laws, for which central differences are singular at shock interfaces. Our Godunov-form flux~\eqref{eq:godunov} replaces them, remaining well-defined across arbitrary jumps.

\paragraph{Flow matching for {PDE}s.}
Stochastic interpolants \citep{albergo2025interpolants} and their function-space lifts \citep{bhola2025fmo} unify flow- and diffusion-based solvers. These methods engineer the generative dynamics; we retain the {VE-SDE} backbone and manipulate the \emph{intrinsic} structure of its score field---the Burgers identity~\eqref{eq:score-burgers} that persists independently of the chosen transport path.

\paragraph{The {PDE} perspective on diffusion.}
\citet{sarkar2026scoreshocks} established that the score field of a {VE-SDE} satisfies a viscous Burgers equation, derived the exact $\tanh$ shock profile (Proposition~5.4), and identified an $\exp(\amp T_{d})$ error-amplification factor (Theorem~6.3) that is sharp under standard parameterizations. \citet{armegioiu2026chaotic} bounded Wasserstein stability for chaotic {PDE} measure transport via average-strain exponents. \textbf{Where we differ:} we convert the score--Burgers identity from an analytical observation into both an architectural prior~\eqref{eq:bv-aware} and a convergence rate~\eqref{eq:thm-improved} that removes the $\exp(\amp T_{d})$ exponent.

\paragraph{Hyperbolic conservation laws.}
Classical solvers---{WENO} \citep{shu1999weno}, discontinuous Galerkin \citep{cockburn2001dg}, and Godunov-type schemes---serve here as ground-truth generators. The Kruzhkov \citep{kruzhkov1970} and Kuznetsov \citep{kuznetsov1976} theories provide the $L^{1}$-contractive semigroup and the vanishing-viscosity rate that scaffold Theorems~\ref{thm:stability}--\ref{thm:improved-rate}.

% ----- Background -------------------------------------------------------
\subsection{Background}
\label{sec:background}

\paragraph{{VE-SDE} preliminaries.}
We adopt the {EDM} formulation \citep{karras2022edm}. Forward: inject Gaussian noise of variance $\sigma^{2}(\tau)$ at diffusion time $\tau\in[0,T_{d}]$, yielding noised density $\ptau=\rhotrue\ast\Gtau$ with score $\score(u,\tau)=\nabla_{u}\log\ptau(u)$. The forward evolution obeys $\dtau\ptau=\Lap_{u}\ptau$, from which~\eqref{eq:score-burgers} follows via Cole--Hopf (Theorem~\ref{thm:double-burgers}). The reverse {SDE} $du = -2\dot\sigma^{2}(\tau)\,\score(u,\tau)\,d\tau + \sqrt{2\dot\sigma^{2}(\tau)}\,d\bar W_{\tau}$ is trained by minimising $\Ldsm$ of $\sth$. The final sample at $\tau=0$ is $\uth$, with $\Law(\uth)$ its distribution.

\paragraph{Hyperbolic target.}
The target scalar conservation law is
\begin{equation}
\dt\mathbf{u}+\dx\flux(\mathbf{u})=0, \label{eq:scalar-law}
\end{equation}
with $C^{1}$ flux (canonical: Burgers, $\flux(u)=u^{2}/2$). Solutions may break down in finite time; the unique Kruzhkov entropy solution $\physsol$ \citep{kruzhkov1970} is selected by the entropy condition, and its semigroup is $L^{1}$-contractive---the analytical backbone of \S\ref{sec:theory}.

\paragraph{Metrics.}
Distributional error is measured by $\Wass{1}$ (Kantorovich--Rubinstein dual); solution regularity is controlled by $\TV(\cdot)$. The compact embedding $\BV(\Omega)\hookrightarrow L^{1}(\Omega)$ (Helly) and the equivalence of $L^{1}$ and $\Wass{1}$ on bounded-{TV} sets provide the compactness used in Theorem~\ref{thm:improved-rate}.

%==============================================================================
% Section 3 · Preliminaries (~0.5 page)
%
% 三块前置知识，做最小集合（每块 ~3-5 行 + 一两个 displayed equation）：
%   §3.1 VE-SDE diffusion model & score matching
%   §3.2 Hyperbolic conservation laws & Kruzhkov entropy solution
%   §3.3 Wasserstein distance & BV space
%==============================================================================
\section{Preliminaries}
\label{sec:prelim}

% ----- §3.1 Diffusion -----
\subsection{Diffusion models and score}
\label{sec:prelim:diffusion}

\todo{State VE-SDE forward process and the score function $\score_{\tau}(u) = \nabla_{u} \log \ptau(u)$ where $\ptau = \rhotrue \ast \Gtau$. Cite Song et al. and EDM \citep{karras2022edm}. Give the reverse-time SDE in one displayed equation:}

\begin{equation}
\label{eq:reverse-sde}
du = -2\,\dot\sigma^{2}(\tau)\,\score_{\tau}(u)\,d\tau + \sqrt{2\,\dot\sigma^{2}(\tau)}\,d\bar W_{\tau}.
\end{equation}

\todo{Mention denoising score matching loss $\Ldsm$ in one line.}

% ----- §3.2 Hyperbolic PDE -----
\subsection{Hyperbolic conservation laws and entropy solutions}
\label{sec:prelim:pde}

\todo{Give the scalar conservation law:}

\begin{equation}
\label{eq:scalar-law}
\dt u + \dx \flux(u) = 0, \qquad u(\cdot, 0) = u_{0},
\end{equation}

\todo{state that classical solutions may fail to exist due to shocks, weak solutions are non-unique, and that uniqueness is restored by the Kruzhkov entropy condition \citep{kruzhkov1970}. State Rankine–Hugoniot and Lax conditions briefly.}

% ----- §3.3 W and BV -----
\subsection{Wasserstein distance and bounded variation}
\label{sec:prelim:wbv}

\todo{Define $\Wass{1}(\mu,\nu)$ via Kantorovich–Rubinstein duality and $\BV$ space. Note the $L^{1}$-contraction property of entropy solutions \citep{kruzhkov1970}, which is the cornerstone of the convergence proofs in Section~\ref{sec:theory}.}

%==============================================================================
% Section 4 · Double-Burgers Structure (~1.5 page) ⭐ 核心观察
%
% 这一节是路径 A 的"故事中点"——把 Score Shocks 的 Burgers 等式与目标 PDE 的
% Burgers 结构耦合起来，呈现"双 Burgers"图景，并给出 Theorem 1。
%
% 子节：
%   §4.1 Score-level Burgers (recap of Score Shocks)
%   §4.2 Solution-level Burgers (target PDE)
%   §4.3 Coupling: shock set correspondence + Theorem 1
%==============================================================================
\section{The Double-Burgers Structure}
\label{sec:double-burgers}

% ----- §4.1 Score-level -----
\subsection{Score-level Burgers}
\label{sec:db:score-level}

\todo{Recap the Score Shocks identity in one displayed equation:}

\begin{equation}
\label{eq:score-burgers}
\dtau \score_{\tau} + 2\,\score_{\tau}\cdot\nabla_{u}\score_{\tau} = \Lap_{u}\score_{\tau},
\end{equation}

\todo{namely $u := -2\score$ satisfies a viscous Burgers equation in $\tau$. Cite \citep{sarkar2026scoreshocks} and emphasize the resulting tanh-shaped interfacial profile of the score near data modes.}

% ----- §4.2 Solution-level -----
\subsection{Solution-level Burgers}
\label{sec:db:soln-level}

\todo{Take the target scalar law \eqref{eq:scalar-law} with $\flux(u) = u^{2}/2$ as the canonical Burgers case. Define the data distribution $\rho_{t}$ at physical time $t$ as the push-forward of the initial-condition distribution $\nu_{0}$ under the PDE flow.}

% ----- §4.3 Coupling + Theorem 1 -----
\subsection{Coupling: physical and score shock sets coincide}
\label{sec:db:coupling}

\todo{Geometric picture: the physical shock set $\Shockphys(t) \subset \Omega$ becomes a singular support of $\rho_{t}$, which determines the location of the score-level interfacial layer $\Shockscore(\tau)$. The two Burgers structures are coupled.}

\begin{theorem}[Double-Burgers coupling]
\label{thm:double-burgers}
\todo{State the formal statement (see path\_A\_method\_skeleton §4 Theorem 1). Sketch: the joint density $p_{\tau, t}(u) = \rho_{t} \ast \Gtau(u)$ satisfies a coupled system in $(\tau, t)$, with the score in $\tau$ direction obeying \eqref{eq:score-burgers} and the data in $t$ direction obeying the Liouville equation along characteristics.}
\end{theorem}

\todo{Add a 1-paragraph proof sketch here; full proof in Appendix~\ref{appx:proofs}.}

% ----- Figure placeholder -----
\begin{figure}[t]
    \centering
    \fbox{\rule[-.5cm]{0cm}{4cm} \rule[-.5cm]{8cm}{0cm}}
    \caption{\todo{Figure 1: schematic of the double-Burgers coupling. Left: physical solution develops a shock at $x_{s}(t)$. Right: score field develops an interfacial tanh layer at $x_{s}(\tau)$ in diffusion time. Key take-away: the two shock sets are coupled, and our method exploits this geometric correspondence to embed shock structure directly into the score architecture.}}
    \label{fig:double-burgers-schematic}
\end{figure}

%==============================================================================
% Section 5 · Method: EntroDiff (~2 pages)
%
% 三个组件 + 损失家族（与 path_A_method_skeleton §3 完全对齐）：
%   §5.1 Viscosity-matched noise schedule
%   §5.2 BV-aware score parameterization (核心 novelty)
%   §5.3 Loss family
%   §5.4 Reverse-time sampler with Godunov-form guidance
%==============================================================================
\section{Method: EntroDiff}
\label{sec:method}

% ----- §5.1 Viscosity-matched schedule -----
\subsection{Viscosity-matched noise schedule}
\label{sec:method:schedule}

\todo{Motivation: from \eqref{eq:score-burgers} the effective viscosity of the score-level Burgers is $g(\tau)^{2}/2$. Aligning it with the physical viscosity of the target PDE makes the two Burgers share the same viscous profile. We choose}

\begin{equation}
\label{eq:visc-matched}
\sigma^{2}(\tau) = 2\,\physvis\,\tau, \qquad \tau \in [0, T_{d}].
\end{equation}

% ----- §5.2 BV-aware parameterization -----
\subsection{BV-aware score parameterization}
\label{sec:method:bv-aware}

\todo{Key novelty. Score Shocks Proposition 5.4 gives the exact tanh form of $\score$ in the interfacial layer. We embed this form directly into the network architecture, decomposing the score into a smooth background gradient plus a tanh-shaped jump:}

\begin{equation}
\label{eq:bv-aware}
\sth(u, \tau) = \nabla \phi^{\mathrm{sm}}_{\theta}(u, \tau)
              + \frac{\kappa_{\theta}(u, \tau)}{2} \tanh\!\left( \frac{\phi^{\mathrm{sh}}_{\theta}(u, \tau)}{2} \right) \nabla \phi^{\mathrm{sh}}_{\theta}(u, \tau),
\end{equation}

\todo{where $\phi^{\mathrm{sm}}$ is a smooth background potential, $\phi^{\mathrm{sh}}$ is a signed-distance-to-shock function (zero on the shock set), and $\kappa$ is the local jump magnitude (constrained by Rankine–Hugoniot). This parameterization severs the $\exp(\amp)$ amplification source identified in Score Shocks Theorem 6.3.}

% ----- §5.3 Loss family -----
\subsection{Loss family}
\label{sec:method:loss}

\todo{Define the four loss components (Table~\ref{tab:loss-family}) and the final composite loss:}

\begin{equation}
\label{eq:loss-final}
\mathcal{L} = \Ldsm + \lambda_{\mathrm{ent}}\,\Rent(\Dth) + \lambda_{\mathrm{BV}}\,\TV(\Dth)
            + \lambda_{\mathrm{Burg}}\,\bigl\| \dtau \sth + 2\,\sth\,\dx \sth - \partial_{uu} \sth \bigr\|^{2}.
\end{equation}

\begin{table}[t]
    \caption{\todo{Loss-family ablation table. Key take-away: each regularizer addresses a specific theoretical role — entropy regularizer enforces Kruzhkov admissibility, BV regularizer keeps $\hat u$ in the entropy-solution function class, Burgers-consistency regularizer pushes $\sth$ onto the score Burgers manifold.}}
    \label{tab:loss-family}
    \centering
    \begin{tabular}{llll}
        \toprule
        Loss & Form & Theoretical role & Empirical role \\
        \midrule
        $\Ldsm$ & DSM & baseline regression target & generalization \\
        $\Lent$ & Kruzhkov entropy regularizer & physical-solution selector & shock-stability \\
        $\Lbv$ & total-variation penalty & keeps $\hat u \in \BV$ & oscillation control \\
        $\Lburg$ & score-Burgers consistency & manifold of admissible scores & robustness near shock \\
        \bottomrule
    \end{tabular}
\end{table}

% ----- §5.4 Sampler -----
\subsection{Reverse-time sampler with Godunov-form guidance}
\label{sec:method:sampler}

\todo{Use EDM-style probability-flow ODE plus DPS-style guidance, but replace the central-difference PDE residual (used by DiffusionPDE) with a Godunov-form entropy-consistent flux difference. This single change is what makes the guidance term well-defined across shocks.}

\begin{equation}
\label{eq:reverse-ode}
\frac{du}{d\tau} = -\sigtau\,\dot\sigma(\tau)\,\sth(u, \tau)
                  - \zeta_{\mathrm{obs}}\,\nabla_{u}\,\mathcal{L}_{\mathrm{obs}}(u)
                  - \zeta_{\mathrm{PDE}}\,\nabla_{u}\,\Lpde^{\mathrm{Godunov}}(u).
\end{equation}

%==============================================================================
% Section 6 · Theory (~2 pages) ⭐ 论文理论核心
%
% NeurIPS 2026 W2 重构：原 04_double_burgers 的 Theorem 1 已并入 §6.1。
%
% 5 个 Theorem，按 path_A_method_skeleton §4：
%   §6.1 Double-Burgers coupling (Theorem 1)        ⭐ 桥段（本步实写）
%   §6.2 Entropy-solution stability (Theorem 2, baseline rate)
%   §6.3 Improved rate via BV-aware (Theorem 3) ⭐ 论文主定理
%   §6.4 Shock-location admissibility (Theorem 4)
%   §6.5 JKO correspondence (Theorem 5)
%
% 每个 Theorem 在正文给 1 段思想（statement + 一句证明思想）；完整证明在
% Appendix~\ref{appx:proofs}。
%==============================================================================
\section{Theory}
\label{sec:theory}

The theorems of this section sequentially erect the convergence theory of \emph{EntroDiff}. §\ref{sec:theory:dburg} establishes the structural foundation by coupling the score-level and solution-level Burgers structures (Theorem~\ref{thm:double-burgers}); §\ref{sec:theory:stability} gives a baseline stability bound that exposes the $\exp(\amp T)$ amplification factor of standard parameterizations; §\ref{sec:theory:improved-rate} states the main contribution---an improved $\Wass{1} \le \mathcal{O}(\varepsilon^{1/2})$ rate that excises this exponent (Theorem~\ref{thm:improved-rate}); §\ref{sec:theory:shock-location} certifies shock-location admissibility via Rankine--Hugoniot and Lax conditions; and §\ref{sec:theory:jko} closes the loop by recasting the reverse process as a {JKO} gradient flow.

% ============================================================================
% §6.1 · Theorem 1 · Double-Burgers coupling (合并自 W2 重构前的 §4)
% ============================================================================
\subsection{Double-Burgers coupling}
\label{sec:theory:dburg}

The point of departure is a structural identity established in \citet{sarkar2026scoreshocks}: along the diffusion-time axis $\tau$, the score field of any VE-SDE diffusion model satisfies a viscous Burgers equation,
\begin{equation}
  \label{eq:score-burgers}
  \dtau \score(u, \tau) + 2\,\score(u, \tau) \cdot \nabla_u \score(u, \tau) = \Lap_u \score(u, \tau).
\end{equation}
Equivalently, $u \mapsto -2\,\score$ is a Burgers velocity field in $(u, \tau)$. When the target distribution $\rhotrue$ is itself supported on solutions of a hyperbolic conservation law of the form~\eqref{eq:scalar-law}, a second Burgers structure emerges in the orthogonal direction: the data law at physical time $t$, $\rho_{t} \coloneqq (\Phi^{\flux}_{t})_{\#} \rho_{0}$ with $\Phi^{\flux}_{t}$ the entropy semigroup of $\flux$, develops shocks at locations $\Shockphys(t) \subset \Omega$ in finite time. The two structures are not merely parallel; they are \emph{coupled}, in that physical shocks define the singular support of $\rho_{t}$ and thereby fix the location of the interfacial layer $\Shockscore(\tau)$ at which~\eqref{eq:score-burgers} develops a $\tanh$-shaped jump.

\begin{theorem}[Double-Burgers coupling]
\label{thm:double-burgers}
Let $\rho_{t}$ denote the law of the Kruzhkov entropy solution of~\eqref{eq:scalar-law} at physical time $t \in [0, T_{\mathrm{phys}}]$, and write $p_{\tau, t} \coloneqq \rho_{t} \ast \Gtau$ for the corresponding noised density at diffusion time $\tau$. Setting $s_{\tau, t} \coloneqq \nabla_u \log p_{\tau, t}$, the joint density $p_{\tau, t}$ satisfies the coupled system
\begin{align}
  \label{eq:dburg-system}
  \dtau s_{\tau, t} &= \Lap_u s_{\tau, t} + 2\,s_{\tau, t} \cdot \nabla_u s_{\tau, t}, \\
  \dt \rho_{t} + \dx\!\bigl(\flux(\mathbf{u})\,\rho_{t}\bigr) &= 0 \qquad \text{(in the entropy-distribution sense)}. \notag
\end{align}
Moreover, for $\tau$ sufficiently small, the interfacial layer $\Shockscore(\tau)$ of $s_{\tau, t}$ is supported in an $\mathcal{O}(\sigma(\tau))$ neighbourhood of the lift of the physical shock set $\Shockphys(t)$ to state space.
\end{theorem}

\todo{Proof sketch (1 paragraph): the $\tau$-direction equation follows from \citet{sarkar2026scoreshocks}, Theorem~4.3, applied pointwise in $t$; the $t$-direction equation is the Liouville continuity equation along the characteristic flow of $\flux$, lifted to the entropy-solution sense via $L^{1}$-contractivity \citep{kruzhkov1970}. The geometric correspondence between $\Shockphys$ and $\Shockscore$ is obtained from the singular support of $\rho_{t}$ at physical shocks combined with Laplace's method on the convolution $\rho_{t} \ast \Gtau$. Full proof in Appendix~\ref{appx:proofs}.}

% ============================================================================
% §6.2 · Theorem 2 · Entropy-solution stability (baseline rate)
% ============================================================================
\subsection{Entropy-solution stability (baseline rate)}
\label{sec:theory:stability}

\begin{theorem}[Entropy-solution stability]
\label{thm:stability}
Let the score training error satisfy $\E_{\tau}\|\sth - \score_{\tau}\|^{2} \le \varepsilon^{2}$. Let $\hat u$ be the sample produced by the reverse sampler under the standard parameterization, and $u^{\star}$ the Kruzhkov entropy solution of \eqref{eq:scalar-law}. Then
\begin{equation}
\label{eq:thm-stability}
\Wass{1}\bigl(\Law(\hat u), \Law(u^{\star})\bigr) \le C_{1}\,\varepsilon\,\exp(\amp T),
\end{equation}
where $\amp = \sup_{\tau} \snr(\tau)/2$ is the Score Shocks amplification exponent.
\end{theorem}

\todo{Proof sketch: Gronwall + Score Shocks Theorem 6.3 + Kruzhkov $L^{1}$-contraction. Full proof in Appendix~\ref{appx:proofs}.}

% ============================================================================
% §6.3 · Theorem 3 · Improved rate (main result)                        ⭐
% ============================================================================
\subsection{Improved rate via BV-aware parameterization}
\label{sec:theory:improved-rate}

\begin{theorem}[Improved rate via BV-aware parameterization]
\label{thm:improved-rate}
Under parameterization \eqref{eq:bv-aware} and loss \eqref{eq:loss-final} with $\lambda_{\mathrm{BV}} > 0$, and under suitable regularity assumptions on the data distribution, we have
\begin{equation}
\label{eq:thm-improved}
\Wass{1}\bigl(\Law(\hat u), \Law(u^{\star})\bigr) \le C_{2}\,\varepsilon^{1/2},
\end{equation}
removing the $\exp(\amp T)$ blow-up factor of Theorem~\ref{thm:stability}.
\end{theorem}

\todo{Proof sketch: parameterization \eqref{eq:bv-aware} inherits the exact tanh interfacial profile from \citet{sarkar2026scoreshocks}, Proposition~5.4, so the Gronwall constant degenerates from $\amp$ to $\mathcal O(1)$. The BV penalty $\lambda_{\mathrm{BV}} \TV(\Dth)$ controls the total variation of $\hat u$, placing it in the $L^{1}$-compact entropy-solution class. Full proof in Appendix~\ref{appx:proofs}.}

\begin{remark}
\label{rem:tightness}
\todo{Comment that the rate $\varepsilon^{1/2}$ is consistent with optimal rates for entropy-solution approximation by viscosity methods \citep{kuznetsov1976}. Whether $\varepsilon^{1}$ is achievable under stronger structural assumptions is an open question.}
\end{remark}

% ============================================================================
% §6.4 · Theorem 4 · Shock-location admissibility
% ============================================================================
\subsection{Shock-location admissibility}
\label{sec:theory:shock-location}

\begin{theorem}[Shock-location admissibility]
\label{thm:shock-location}
In the limit $\sigma(\tau) \to 0$, the shock $x_{s}(t)$ of $\hat u$ satisfies the Rankine--Hugoniot condition
\begin{equation}
\label{eq:rh}
\dot x_{s}(t) = \frac{\flux(u_{L}) - \flux(u_{R})}{u_{L} - u_{R}},
\end{equation}
and the Lax entropy condition $\flux'(u_{L}) \ge \dot x_{s} \ge \flux'(u_{R})$ holds almost everywhere.
\end{theorem}

\todo{Proof sketch: Score Shocks Theorem 5.11 (speciation time) + viscosity-matched schedule \eqref{eq:visc-matched} aligns score-level shock formation with physical-time shock formation. R--H is then inherited from the score-level jump correspondence (Score Shocks Theorem 5.5). Full proof in Appendix~\ref{appx:proofs}.}

% ============================================================================
% §6.5 · Theorem 5 · JKO correspondence
% ============================================================================
\subsection{{JKO} correspondence}
\label{sec:theory:jko}

\begin{theorem}[{JKO} correspondence]
\label{thm:jko}
Under viscosity-matched schedule \eqref{eq:visc-matched} and loss term $\Lburg$, the reverse-time ODE \eqref{eq:reverse-ode} in the time-discretization limit $\Delta\tau \to 0$ corresponds to the {JKO} step
\begin{equation}
\label{eq:jko}
\rho^{n+1} = \arg\min_{\rho \in \mathcal{P}_{2}}
            \left\{ \mathcal{F}(\rho) + \frac{1}{2\Delta\tau}\,\Wass{2}^{2}(\rho, \rho^{n}) \right\},
\end{equation}
where $\mathcal{F}(\rho) = \Ent(\rho \,|\, \rhotrue) + \langle \text{PDE constraint} \rangle$.
\end{theorem}

\todo{Proof sketch: classical Jordan--Kinderlehrer--Otto \citep{jko1998} gives the Wasserstein gradient-flow structure of Fokker--Planck; the PDE constraint enters as a Lagrangian multiplier. Full proof in Appendix~\ref{appx:proofs}.}

\input{sections/07_experiments.tex}
\input{sections/08_conclusion.tex}
%------------------------------------------------------------------------------

\begin{ack}
% 双盲提交时此段被自动隐藏
\end{ack}

%------------------------------------------------------------------------------
% References
%------------------------------------------------------------------------------
\bibliographystyle{plainnat}
\bibliography{references}

%==============================================================================
\appendix
%==============================================================================
% Appendix A1 · Complete Proofs of Theorems 1–5
%
% 每个 Theorem 的 proof sketch 在 §6 main body 中；
% 完整 ε-δ 推导在此附录。
%==============================================================================
\section{Proofs of Main Theorems}
\label{appx:proofs}

% ----- A1.1 · Theorem 1 · Double-Burgers coupling -----
\subsection{Proof of Theorem~\ref{thm:double-burgers} (Double-Burgers coupling)}
\label{appx:proofs:thm1}
% Theorem 1 (Double-Burgers Coupling) -- Appendix Proof
% Included under \subsection{Proof of Theorem~\ref{thm:double-burgers}} in A1_proofs.tex

\begin{theorem}[Double-Burgers coupling --- restated]
\label{thm:double-burgers-app}
Let $\Omega \subset \R$ be the physical domain and consider the 1D scalar conservation law
\begin{equation}
\dt u + \dx \flux(u) = 0, \qquad \flux \in C^{2}(\R), \quad u(0,\cdot) = \initdat .
\label{eq:thm1-claw}
\end{equation}
Let $S_{t}$ denote the Kruzhkov entropy semigroup \citep{kruzhkov1970}. Let the initial data be random: $\initdat \sim \rho_{0}$ with $\rho_{0} \in \mathcal{P}(L^{\infty}(\Omega) \cap \BV(\Omega))$, and define the pushforward distribution at physical time $t$:
\begin{equation}
\rhotphys \coloneqq (S_{t})_{\#} \rho_{0}.
\label{eq:thm1-pushforward}
\end{equation}
At each fixed $t$, apply Gaussian smoothing with diffusion time $\tau > 0$:
\begin{equation}
p_{\tau, t}(u) \coloneqq (\rhotphys * \Gtau)(u), \qquad
\Gtau(u) = (4\pi\tau)^{-d/2} \exp\!\left(-\frac{\norm{u}^{2}}{4\tau}\right),
\label{eq:thm1-gaussian-smoothing}
\end{equation}
where $u \in \R^{d}$ is the state variable (for scalar conservation laws, $d = 1$; the Score Burgers part holds for all $d \ge 1$). Define the score and its Cole--Hopf conjugate:
\begin{equation}
\score_{\tau, t}(u) \coloneqq \nabla_{u} \log p_{\tau, t}(u), \qquad
w_{\tau, t}(u) \coloneqq -2\, \score_{\tau, t}(u).
\label{eq:thm1-score-def}
\end{equation}
Then the following hold:
\begin{enumerate}
\item[(A)] \textbf{Score Burgers ($\tau$-direction).} For each fixed $t$, for all $\tau > 0$ where $p_{\tau, t} > 0$,
\begin{equation}
\boxed{\; \dtau w_{\tau, t} + (w_{\tau, t} \cdot \nabla_{u}) w_{\tau, t} = \Lap_{u} w_{\tau, t} \;}
\label{eq:thm1-score-burgers}
\end{equation}
with $\nabla \times w_{\tau, t} = 0$ identically (since $w = -2\nabla \log p$).

\item[(B)] \textbf{Physical transport ($t$-direction).} The curve $t \mapsto \rhotphys$ in $\mathcal{P}_{2}(\R^{d})$ satisfies the continuity equation
\begin{equation}
\boxed{\; \dt \rhotphys + \nabla_{u} \cdot (\rhotphys \, V_{t}) = 0, \qquad
V_{t}(u)(x) = -\dx \flux(u(x)) \;}
\label{eq:thm1-physical-transport}
\end{equation}
in the sense of distributions, obtained as the vanishing-viscosity limit $\varepsilon \to 0^{+}$ of the regularized evolution.

\item[(C)] \textbf{Shock set co-location.} Let
\begin{equation}
\Shockphys(t) \coloneqq \{\, x \in \Omega : u(\cdot, t) \text{ has a jump at } x \,\}
\label{eq:thm1-phys-shock}
\end{equation}
be the physical shock set of the entropy solution. For each $\tau > 0$, define the score interfacial layer set
\begin{equation}
\Shockscore(\tau, t) \coloneqq \{\, u \in \R^{d} : \norm{\nabla_{u} w_{\tau, t}(u)} \gtrsim 1 / \tau \,\},
\label{eq:thm1-score-shock}
\end{equation}
i.e., the region where the score gradient is large (the $\tanh$-layer identified by Score Shocks \citep{sarkar2026scoreshocks}). Let $\pi_{\Omega}: \R^{d} \to \Omega$ be the projection from state space to physical space (identifying the spatial coordinate of the shock value). Then
\begin{equation}
\boxed{\; \pi_{\Omega}\big(\Shockscore(\tau, t)\big) \;\xrightarrow{\; \tau \to 0^{+} \;}\; \Shockphys(t) \qquad \text{in Hausdorff distance.} \;}
\label{eq:thm1-shock-colocation}
\end{equation}
\end{enumerate}
\end{theorem}

\paragraph{Assumptions.}
\begin{enumerate}
\item[(A1)] \textbf{Gaussian smoothing.} For each fixed $t$, $\rhotphys$ is a finite Borel measure, and the convolution $p_{\tau, t} = \rhotphys * \Gtau$ is strictly positive and $C^{\infty}$ for all $\tau > 0$.

\item[(A2)] \textbf{Entropy semigroup.} $S_{t}$ is the Kruzhkov entropy semigroup, so $\rhotphys = (S_{t})_{\#} \rho_{0}$ is well-defined and $t \mapsto \rhotphys$ is continuous in the narrow topology.

\item[(A3)] \textbf{BV regularity.} $\rho_{0}$ is supported on $\BV(\Omega) \cap L^{\infty}(\Omega)$. Consequently, for each $t$, $u(\cdot, t) \in \BV(\Omega)$ with $\TV(u(\cdot, t)) \le \TV(\initdat)$, and the jump set $\Shockphys(t)$ is a finite (or countably infinite) set of rectifiable curves in $(t,x)$-space.

\item[(A4)] \textbf{Convex flux, for concreteness.} $\flux$ is uniformly convex ($\flux'' \ge c > 0$); the non-convex case follows via Oleinik's generalization \citep{oleinik1957} and is discussed in the appendix.
\end{enumerate}

\begin{proof}
We establish the three structural identities (A), (B), (C) in order.

\medskip
\noindent
\textbf{Part A --- Score Burgers ($\tau$-direction).}
This part follows directly from the Cole--Hopf correspondence and is essentially a restatement of Score Shocks Theorem~4.3 \citep{sarkar2026scoreshocks}.

By definition, $p_{\tau, t} = \rhotphys * \Gtau$. Since $\Gtau$ is the heat kernel, $\dtau \Gtau = \Lap_{u} \Gtau$, and differentiation under the integral (justified by A1) yields
\begin{equation}
\dtau p_{\tau, t} = \Lap_{u} p_{\tau, t}.
\label{eq:thm1-A1}
\end{equation}

Set $\varphi \coloneqq \log p_{\tau, t}$ ($p_{\tau, t} > 0$ by A1). Then
\begin{align}
\dtau \varphi
&= \frac{\dtau p}{p}
 = \frac{\Lap p}{p}
 = \Lap \varphi + \abs{\nabla \varphi}^{2},
\label{eq:thm1-A2}
\end{align}
where we used $\nabla \cdot (p \nabla \varphi) = \nabla p \cdot \nabla \varphi + p \Lap \varphi = p \abs{\nabla \varphi}^{2} + p \Lap \varphi$ and $\nabla \cdot (p \nabla \varphi) = \nabla \cdot (\nabla p) = \Lap p$, so $\Lap p / p = \Lap \varphi + \abs{\nabla \varphi}^{2}$.

Now $\score_{\tau, t} = \nabla_{u} \varphi$. Taking the gradient of~\eqref{eq:thm1-A2},
\begin{align}
\dtau \score_{\tau, t}
&= \nabla(\Lap \varphi) + \nabla(\abs{\nabla \varphi}^{2})
 = \Lap \score_{\tau, t} + \nabla(\abs{\score_{\tau, t}}^{2}).
\label{eq:thm1-A3}
\end{align}

Since $\score_{\tau, t} = \nabla \varphi$ is a gradient, $\nabla \score_{\tau, t} = \nabla^{2} \varphi$ is symmetric: $\partial_{u_{j}} s_{i} = \partial_{u_{i}} s_{j}$. Hence
\begin{align}
[\nabla(\abs{\score}^{2})]_{i}
&= \partial_{u_{i}} \sum_{j} s_{j}^{2}
 = 2 \sum_{j} s_{j} \, \partial_{u_{i}} s_{j}
 = 2 \sum_{j} s_{j} \, \partial_{u_{j}} s_{i}
 = 2[(\score \cdot \nabla) \score]_{i}.
\label{eq:thm1-A4}
\end{align}

Substituting~\eqref{eq:thm1-A4} into~\eqref{eq:thm1-A3} gives the score PDE:
\begin{equation}
\dtau \score_{\tau, t} = \Lap_{u} \score_{\tau, t} + 2 (\score_{\tau, t} \cdot \nabla_{u}) \score_{\tau, t}.
\label{eq:thm1-A5}
\end{equation}

Finally, define $w_{\tau, t} \coloneqq -2\, \score_{\tau, t}$. Then
\begin{align}
\dtau w
&= -2\,\dtau \score
 = -2\bigl(\Lap \score + 2(\score \cdot \nabla) \score\bigr)
 = \Lap w - 2(w \cdot \nabla) \score.
\label{eq:thm1-A6}
\end{align}

Since $\score = -w/2$, we have $(w \cdot \nabla)\score = -\tfrac{1}{2} (w \cdot \nabla) w$, hence
\begin{equation}
\dtau w = \Lap w + (w \cdot \nabla) w.
\label{eq:thm1-A7}
\end{equation}

Rearranging yields the viscous Burgers equation~\eqref{eq:thm1-score-burgers}. The curl-free condition $\nabla \times w = -\nabla \times (2\nabla \log p) = 0$ is automatic.

\medskip
\noindent
\textbf{Part B --- Physical Transport ($t$-direction).}
We prove~\eqref{eq:thm1-physical-transport} via the vanishing viscosity method, following the structure sketched in \citet{kruzhkov1970} and adapted to the distributional setting.

\paragraph{B.1 Viscous regularization.}
Introduce the viscous perturbation of the conservation law:
\begin{equation}
\dt u^{\varepsilon} + \dx \flux(u^{\varepsilon}) = \varepsilon\, \dx^{2} u^{\varepsilon}, \qquad
u^{\varepsilon}(0, \cdot) = \initdat .
\label{eq:thm1-B1}
\end{equation}

Let $S_{t}^{\varepsilon}$ be the solution operator of~\eqref{eq:thm1-B1}. For BV initial data, the solution $u^{\varepsilon}(\cdot, t)$ is smooth for all $t > 0$ (parabolic regularization). Define the regularized pushforward distribution
\begin{equation}
\rho_{t}^{\varepsilon} \coloneqq (S_{t}^{\varepsilon})_{\#} \rho_{0}.
\label{eq:thm1-B1-dist}
\end{equation}

By Kruzhkov's theory \citep{kruzhkov1970}, for any $T < \infty$,
\begin{equation}
u^{\varepsilon}(\cdot, t) \to u(\cdot, t) \quad \text{in } L^{1}_{\mathrm{loc}}(\Omega), \quad \text{uniformly in } t \in [0,T],
\label{eq:thm1-B1-conv}
\end{equation}
as $\varepsilon \to 0^{+}$, where $u$ is the unique entropy solution. In particular, $\rho_{t}^{\varepsilon} \rightharpoonup \rhotphys$ in the narrow topology.

\paragraph{B.2 Liouville equation for the viscous flow.}
For the viscous PDE~\eqref{eq:thm1-B1}, the solution at each fixed $x$ evolves according to
\begin{equation}
\dt u^{\varepsilon}(t, x) = -\dx \flux(u^{\varepsilon}(t, x)) + \varepsilon\, \dx^{2} u^{\varepsilon}(t, x).
\label{eq:thm1-B2-ode}
\end{equation}

Treating this as an ODE in the state variable $u$ for each fixed $x$, the velocity field is
\begin{equation}
V_{t}^{\varepsilon}(u)(x) \coloneqq -\dx \flux(u) + \varepsilon\, \dx^{2} u.
\label{eq:thm1-B2-V}
\end{equation}

Let $\psi \in C_{c}^{\infty}(\R^{d})$ be a test function in state space. By the chain rule and the definition $\rho_{t}^{\varepsilon} = (S_{t}^{\varepsilon})_{\#} \rho_{0}$,
\begin{align}
\frac{d}{dt} \int_{\R^{d}} \psi(u) \, d\rho_{t}^{\varepsilon}(u)
&= \E_{\initdat \sim \rho_{0}}\!\bigl[ \nabla \psi(u^{\varepsilon}(t, \cdot)) \cdot \dt u^{\varepsilon}(t, \cdot) \bigr] \notag \\
&= \int_{\R^{d}} \nabla \psi(u) \cdot V_{t}^{\varepsilon}(u) \, d\rho_{t}^{\varepsilon}(u)
   + \varepsilon \int_{\R^{d}} \Lap \psi(u) \, d\rho_{t}^{\varepsilon}(u),
\label{eq:thm1-B2-chain}
\end{align}
where we used $\E[\nabla \psi \cdot \varepsilon \dx^{2} u] = -\varepsilon \E[\nabla^{2} \psi : \dx u \otimes \dx u + \Lap \psi]$ after an integration by parts in $x$; the second-order term simplifies to $\varepsilon \Lap_{u} \psi$ in the state-space measure.

Integration by parts in $u$ yields the weak form
\begin{equation}
\int \psi \, \dt \rho_{t}^{\varepsilon} \, du
= - \int \nabla \psi \cdot V_{t}^{\varepsilon} \; \rho_{t}^{\varepsilon} \, du
+ \varepsilon \int \Lap \psi \; \rho_{t}^{\varepsilon} \, du,
\label{eq:thm1-B2-weak}
\end{equation}
which is the distributional form of
\begin{equation}
\dt \rho_{t}^{\varepsilon} + \nabla_{u} \cdot (\rho_{t}^{\varepsilon} V_{t}^{\varepsilon}) = \varepsilon \Lap_{u} \rho_{t}^{\varepsilon}.
\label{eq:thm1-B2-liouville}
\end{equation}

\paragraph{B.3 Vanishing viscosity limit.}
Take $\varepsilon \to 0^{+}$ in~\eqref{eq:thm1-B2-liouville}. The second-order term $\varepsilon \Lap_{u} \rho_{t}^{\varepsilon} \to 0$ in the sense of distributions. By the $L^{1}$-contraction property of the entropy semigroup and BV compactness \citep{helly, kruzhkov1970}:
\begin{itemize}
\item $\rho_{t}^{\varepsilon} \rightharpoonup \rhotphys$ narrowly;
\item $V_{t}^{\varepsilon}(u)(x) = -\dx \flux(u) + \varepsilon \dx^{2} u \to -\dx \flux(u) =: V_{t}(u)(x)$ almost everywhere in $(t,x)$.
\end{itemize}

Passing to the limit in the weak formulation of~\eqref{eq:thm1-B2-liouville} gives
\begin{equation}
\dt \rhotphys + \nabla_{u} \cdot (\rhotphys \, V_{t}) = 0, \qquad
V_{t}(u)(x) = -\dx \flux(u(x)),
\label{eq:thm1-B3-limit}
\end{equation}
in the distributional sense. This is precisely equation~\eqref{eq:thm1-physical-transport}.

\medskip
\noindent
\textbf{Part C --- Shock Set Co-location.}
This is the core geometric observation of Theorem~1. We sketch the argument; a fully rigorous measure-theoretic proof (including the Hausdorff convergence rate) appears in the supplementary material.

\paragraph{C.1 BV structure of $\rhotphys$.}
By Assumption~A3, $u(\cdot, t) \in \BV(\Omega)$. The Lebesgue decomposition of $u(\cdot, t)$ splits it into an absolutely continuous part and a jump part:
\begin{equation}
u(\cdot, t) = u^{\mathrm{ac}}(\cdot, t) + \sum_{x_{s} \in \Shockphys(t)} \llbracket u \rrbracket_{x_{s}} \, \mathbf{1}_{[x_{s}, \infty)},
\label{eq:thm1-C1-decomp}
\end{equation}
where $\llbracket u \rrbracket_{x_{s}} \coloneqq u(x_{s}^{+}, t) - u(x_{s}^{-}, t)$ is the jump magnitude at shock $x_{s}$.

Correspondingly, the data distribution decomposes as
\begin{equation}
\rhotphys = \rho_{t}^{\mathrm{ac}} + \rho_{t}^{\mathrm{sing}},
\label{eq:thm1-C1-rho-decomp}
\end{equation}
where $\rho_{t}^{\mathrm{sing}}$ is a sum of weighted Dirac masses (or narrow concentrations) at the shock values $\{u(x_{s}^{-}, t), u(x_{s}^{+}, t)\}_{x_{s} \in \Shockphys(t)}$.

\paragraph{C.2 Mode boundary formation under Gaussian smoothing.}
Consider a single shock at $x_{s}$ with left/right limits $u_{L}, u_{R}$. Locally (in state space near $u_{L}, u_{R}$), $\rhotphys$ behaves approximately as
\begin{equation}
\rhotphys \approx w_{L} \delta_{u_{L}} + w_{R} \delta_{u_{R}} + \text{smooth background},
\label{eq:thm1-C2-approx}
\end{equation}
where $w_{L}, w_{R} > 0$ are the spatial weights on each side of the shock.

Gaussian convolution with $\Gtau$ yields
\begin{equation}
p_{\tau, t}(u) \approx w_{L} \Gtau(u - u_{L}) + w_{R} \Gtau(u - u_{R}) + \text{smooth correction}.
\label{eq:thm1-C2-mixture}
\end{equation}

This is a two-component Gaussian mixture. By Score Shocks Proposition~3.1 (speciation threshold) \citep{sarkar2026scoreshocks}, the score field $\score_{\tau, t} = \nabla_{u} \log p_{\tau, t}$ develops an interfacial layer --- a narrow region where $\abs{\nabla_{u} \score}$ is large --- located near the ``mode boundary,'' i.e., the set of $u$ where the two Gaussian components have equal weight. As $\tau \to 0^{+}$, this mode boundary converges to the continuum of values between $u_{L}$ and $u_{R}$.

\paragraph{C.3 Interfacial profile and shock set convergence.}
By Score Shocks Proposition~5.4 \citep{sarkar2026scoreshocks} (the interfacial profile result), near each mode boundary the score field admits the exact decomposition
\begin{equation}
w_{\tau, t}(u) = w^{\mathrm{sm}}(u) + \frac{\kappa(u)}{2} \tanh\!\left(\frac{\phi(u)}{2}\right) \nabla \phi(u) + o(1),
\label{eq:thm1-C3-decomp}
\end{equation}
where $\phi(u)$ is the signed log-ratio of the local mode densities and $\kappa$ is the interfacial strength (related to the jump magnitude $\abs{u_{L} - u_{R}}$). The interfacial layer set is precisely
\begin{equation}
\Shockscore(\tau, t) = \{\, u : \norm{\nabla_{u} w_{\tau, t}(u)} \gtrsim 1/\tau \,\},
\label{eq:thm1-C3-layer}
\end{equation}
which, to leading order, coincides with the set where $\abs{\nabla \phi}^{-1} \lesssim \sqrt{\tau}$.

Now, the key observation: the mode boundaries of $\rhotphys$ --- and therefore the interfacial layers of $w_{\tau, t}$ --- are in one-to-one correspondence with the shock jumps of $u(\cdot, t)$. This is because:
\begin{enumerate}
\item[1.] \textbf{Each physical shock at $x_{s}$ contributes a sharp bimodal feature to $\rhotphys$} (the Dirac-like concentrations at $u_{L}, u_{R}$).
\item[2.] \textbf{Gaussian smoothing of a bimodal feature always produces an interfacial layer} in the score field --- this is a deterministic consequence of the heat-kernel/Cole--Hopf dynamics (Score Shocks Theorem~5.5).
\item[3.] \textbf{Conversely}, smooth (non-shock) regions of $u(\cdot, t)$ contribute only unimodal or slowly varying features to $\rhotphys$, which produce no interfacial layer.
\end{enumerate}

Therefore, projecting $\Shockscore(\tau, t)$ to physical space via $\pi_{\Omega}$ (which maps each state value $u$ to the spatial location where that value occurs) yields a set that converges to $\Shockphys(t)$ as $\tau \to 0^{+}$.

\paragraph{C.4 Hausdorff convergence.}
Formally, for each shock $x_{s} \in \Shockphys(t)$ with jump values $u_{L}, u_{R}$, the interfacial layer in $w_{\tau, t}$ at diffusion time $\tau$ has width $\ell_{\tau} \sim \sqrt{\tau} / \abs{u_{L} - u_{R}}$ (Score Shocks Proposition~5.4). Its spatial projection satisfies
\begin{equation}
\mathrm{dist}_{H}\!\bigl(\pi_{\Omega}(\Shockscore(\tau, t)), \, \Shockphys(t)\bigr) \le C \, \sqrt{\tau},
\label{eq:thm1-C4-bound}
\end{equation}
where the constant $C$ depends on the shock strength and the local regularity of $u(\cdot, t)$. The proof of this bound relies on the BV structure (to control the number of shocks and their separation) and the $\tanh$-profile asymptotics of Score Shocks. The detailed computation is deferred to the supplementary material.
\end{proof}

Assembling Parts~A, B, and~C, we have established:
\begin{itemize}
\item In the \textbf{diffusion time} $\tau$, the score field evolves as a vector viscous Burgers equation (A) --- this is the Cole--Hopf structure that underlies all diffusion models.
\item In the \textbf{physical time} $t$, the data distribution is transported by the entropy-solution dynamics, satisfying a Liouville equation (B) --- this links the PDE to the distributional level.
\item The \textbf{shock sets} of the physical solution and the score field coincide in the $\tau \to 0^{+}$ limit (C) --- this is the geometric coupling that motivates the EntroDiff architecture: the network's $\tanh$-interfacial-layer structure can be aligned with the known shock locations, enabling the exponential-error-amplification removal of Theorem~3.
\end{itemize}

Theorem~1 is thus the structural foundation on which the entire EntroDiff methodology is built.


% ----- A1.2 · Theorem 2 · Entropy-solution stability -----
\subsection{Proof of Theorem~\ref{thm:stability} (Entropy-solution stability)}
\label{appx:proofs:thm2}
% Theorem 2 (Entropy-Solution Stability -- Baseline Rate) -- Appendix Proof
% Included under \subsection{Proof of Theorem~\ref{thm:stability}} in A1_proofs.tex

\begin{theorem}[Entropy-solution stability --- restated]
\label{thm:stability-app}
Let $\physsol$ be the Kruzhkov entropy solution of the scalar conservation law $\dt \mathbf{u} + \dx \flux(\mathbf{u}) = 0$ at a fixed physical time $T_{\mathrm{phys}}$. View the solution as a state vector $u \in \R^{d}$ ($d = N_{x}$ is the number of spatial grid points) and define the target distribution
\begin{equation}
\rhotrue \coloneqq \Law(\physsol).
\label{eq:thm2-rho-star}
\end{equation}

The forward diffusion process (VE-SDE, schedule $\sigtau$):
\begin{equation}
\ptau \coloneqq \rhotrue \ast \Gtau, \qquad
\Gtau(u) = (4\pi\tau)^{-d/2} \exp\!\left(-\frac{\norm{u}^{2}}{4\tau}\right), \quad \tau \in [0, T_{d}].
\label{eq:thm2-ptau}
\end{equation}

The true score field:
\begin{equation}
\score(u, \tau) \coloneqq \nabla_{u} \log \ptau(u).
\label{eq:thm2-score}
\end{equation}

The reverse ODE (standard EDM form, $\tau$ integrated from $T_{d}$ to $0$):

True trajectory:
\begin{align}
\frac{d\unat}{d\tau} &= b(\unat, \tau), \qquad
b(u, \tau) \coloneqq -\sigtau\,\dot\sigma(\tau)\,\score(u, \tau),
\quad \unat(T_{d}) \sim p_{T_{d}},
\label{eq:thm2-reverse-true}
\end{align}

Network-approximated trajectory:
\begin{align}
\frac{d\uhat}{d\tau} &= \hat b(\uhat, \tau), \qquad
\hat b(u, \tau) \coloneqq -\sigtau\,\dot\sigma(\tau)\,\sth(u, \tau),
\quad \uhat(T_{d}) \sim p_{T_{d}}.
\label{eq:thm2-reverse-approx}
\end{align}

Notation: $\unat(\tau)$ is the true-score-driven reverse trajectory; $\uhat(\tau)$ is the learned-score-driven reverse trajectory. $\norm{\cdot}$ denotes the Euclidean norm on $\R^{d}$. $\Wass{1}(\mu, \nu)$ is the 1-Wasserstein distance (Kantorovich--Rubinstein dual). The expectation $\E$ is taken over randomness (initial noise $+$ learned score approximation).

\medskip
\noindent
\textbf{Main theorem.} Under Assumptions~(A1)--(A5), the generated sample distribution $\Law(\uhat(0))$ and the target distribution $\rhotrue$ satisfy
\begin{equation}
\boxed{\;
\Wass{1}\!\bigl(\Law(\uhat(0)),\; \rhotrue\bigr)
\;\le\; C\,\varepsilon\,\exp(\amp T_{d})
\;},
\label{eq:thm2-main}
\end{equation}
where $\amp \coloneqq \sup_{\tau \in [0, T_{d}]} \abs{\sigtau\,\dot\sigma(\tau)}\,L_{s}(\tau)$, $L_{s}(\tau)$ is the Lipschitz constant of $\nabla_{u} \score(\cdot, \tau)$, and $\varepsilon$ is the score-matching training error bound. The constant $C = C_{\sigma} \sqrt{T_{d}}\,(1 + C_{\mathrm{stab}})$ depends only on the schedule upper bound $C_{\sigma}$ and the distributional stability constant $C_{\mathrm{stab}}$, and is independent of $\varepsilon$ and $\amp$.

\medskip
\noindent
\textbf{Connection to the PDE solution.} Equation~\eqref{eq:thm2-main} provides a quantitative guarantee that the diffusion model's sampled solution distribution is close to the PDE solution distribution. If the training data originates from a numerical solver, the Kruzhkov $L^{1}$-contractivity ensures that the numerical error does not amplify during physical time propagation; the bound~\eqref{eq:thm2-main} transfers stably to the true PDE solution (see \S\,PDE Connection below).
\end{theorem}

\paragraph{Assumptions.}
\begin{enumerate}
\item[(A1)] \textbf{Score gradient Lipschitz regularity.} For each $\tau > 0$, the true score field $\score(\cdot, \tau)$ is $C^{1}$ and its Jacobian's operator norm is uniformly bounded: there exists $L_{s}(\tau) < \infty$ such that
\begin{equation}
\norm{\nabla_{u} \score(u, \tau)}_{\mathrm{op}} \le L_{s}(\tau), \quad \forall\, u \in \R^{d},\; \tau \in (0, T_{d}].
\label{eq:thm2-assum-A1}
\end{equation}
Define $\amp(\tau) \coloneqq \abs{\sigtau\,\dot\sigma(\tau)}\,L_{s}(\tau)$ and set $\amp \coloneqq \sup_{\tau \in [0, T_{d}]} \amp(\tau) < \infty$.

\textit{Verifiability:} Since $\ptau = \rhotrue \ast \Gtau$ and $\Gtau$ is the heat kernel, for any $\tau > 0$, $\ptau \in C^{\infty}$ and $\nabla_{u} \score = \nabla_{u}^{2} \log \ptau$ is uniformly bounded. $\rhotrue$ is a pushforward of a BV measure (containing Dirac point masses at shocks); smoothness holds after convolution. $L_{s}(\tau) \sim 1/\tau$ may diverge as $\tau \to 0$, but we only require $\sup_{\tau \in [\tau_{\mathrm{end}}, T_{d}]} L_{s}(\tau) < \infty$ (early-stopping $\tau_{\mathrm{end}} > 0$), consistent with \citet{sarkar2026scoreshocks}.

\item[(A2)] \textbf{Noise schedule boundedness.} There exists $C_{\sigma} > 0$ such that
\begin{equation}
\abs{\sigtau\,\dot\sigma(\tau)} \le C_{\sigma}, \quad \forall\, \tau \in [0, T_{d}].
\label{eq:thm2-assum-A2}
\end{equation}

\textit{Verifiability:} For VE-type schedules (e.g., $\sigma^{2}(\tau) = 2\physvis \tau$), $\sigma \dot\sigma = \physvis$ is constant, which directly satisfies this.

\item[(A3)] \textbf{Score training error (weighted $L^{2}$ bound).} The learned score network $\sth$ satisfies, under the true data distribution $\ptau$,
\begin{equation}
\int_{0}^{T_{d}} \E_{u \sim \ptau}\!\bigl[\,\norm{\sth(u, \tau) - \score(u, \tau)}^{2}\,\bigr] \, d\tau \;\le\; \varepsilon^{2}.
\label{eq:thm2-assum-A3}
\end{equation}

\textit{Verifiability:} This is the training objective of denoising score matching \citep{vincent2011connection} (with weight $\sigtau^{2}$); $\varepsilon$ is determined by the training loss.

\item[(A4)] \textbf{Distributional stability (Kantorovich--Rubinstein control).} There exists a constant $C_{\mathrm{stab}} > 0$ such that for any $\tau \in [0, T_{d}]$ and any Lipschitz function $f: \R^{d} \to \R$,
\begin{equation}
\abs{\,\E_{\uhat(\tau)}[f(\uhat)] - \E_{\unat(\tau)}[f(\unat)]\,}
\;\le\;
C_{\mathrm{stab}} \cdot \Wass{1}\!\bigl(\Law(\uhat(\tau)),\, \Law(\unat(\tau))\bigr) \cdot \Lip(f),
\label{eq:thm2-assum-A4}
\end{equation}
where $\Lip(f) \coloneqq \sup_{u \neq v} \abs{f(u) - f(v)} / \norm{u - v}$.

\textit{Verifiability:} This is a direct generalization of the Kantorovich--Rubinstein dual $\abs{\int f d\mu - \int f d\nu} \le \Lip(f) \cdot \Wass{1}(\mu, \nu)$, for which $C_{\mathrm{stab}} = 1$. Introducing $C_{\mathrm{stab}}$ accommodates additional smoothness offsets introduced by the score network; under compactly supported measures and smooth score assumptions, $C_{\mathrm{stab}} = \mathcal{O}(1)$.

\item[(A5)] \textbf{Network Lipschitz bound.} The learned score network $\sth(\cdot, \tau)$ is uniformly Lipschitz across all $\tau \in [0, T_{d}]$: there exists $L_{s_{\theta}} < \infty$ independent of $\tau$ such that for all $u, v \in \R^{d}$,
\begin{equation}
\norm{\sth(u, \tau) - \sth(v, \tau)} \le L_{s_{\theta}} \norm{u - v}.
\label{eq:thm2-assum-A5}
\end{equation}

\textit{Verifiability:} Standard MLPs/UNets are Lipschitz on compact input domains; boundedness of weights is ensured by spectral normalization or training regularization.
\end{enumerate}

\begin{proof}
The proof proceeds in nine steps. Steps~1--3 set up the synchronous coupling framework; Steps~4--5 address the core technical difficulty (measure mismatch); Steps~6--8 synthesize the Gronwall estimate; Step~9 yields the final $\Wass{1}$ bound.

\medskip
\noindent
\textbf{Step~1 --- Synchronous coupling.}
Define the error process
\begin{equation}
\errsc(\tau) \coloneqq \uhat(\tau) - \unat(\tau), \qquad \tau \in [0, T_{d}].
\label{eq:thm2-step1}
\end{equation}

Since both trajectories share the same initial distribution $p_{T_{d}}$, we have $\errsc(T_{d}) = 0$ almost surely. Differentiating along the reverse ODE with respect to $\tau$ (with $\tau$ decreasing from $T_{d}$ to $0$):
\begin{equation}
\frac{d}{d\tau} \errsc(\tau)
= \hat b(\uhat(\tau), \tau) - b(\unat(\tau), \tau).
\label{eq:thm2-step1-ode}
\end{equation}

Add-and-subtract decomposition:
\begin{align}
\frac{d}{d\tau} \errsc(\tau)
&= \bigl[b(\uhat(\tau), \tau) - b(\unat(\tau), \tau)\bigr]
   \;+\;
   \bigl[\hat b(\uhat(\tau), \tau) - b(\uhat(\tau), \tau)\bigr] \notag \\[4pt]
&\eqqcolon (\mathrm{I}) + (\mathrm{II}).
\label{eq:thm2-step1-decomp}
\end{align}

The two terms differ in nature:
\begin{itemize}
\item $(\mathrm{I})$: propagation error due to trajectory deviation $\uhat \neq \unat$ under the \emph{same} drift field $b$;
\item $(\mathrm{II})$: model error due to the score approximation $\sth \neq \score$ in the drift field itself.
\end{itemize}

\medskip
\noindent
\textbf{Step~2 --- Controlling $(\mathrm{I})$ (propagation error).}
From $b(u, \tau) = -\sigtau\,\dot\sigma(\tau)\,\score(u, \tau)$ and Assumption~(A1):
\begin{align}
\norm{b(\uhat, \tau) - b(\unat, \tau)}
&\le \abs{\sigtau\,\dot\sigma(\tau)} \cdot \norm{\score(\uhat, \tau) - \score(\unat, \tau)} \notag \\
&\le \abs{\sigtau\,\dot\sigma(\tau)} \cdot L_{s}(\tau) \cdot \norm{\uhat - \unat} \notag \\
&= \amp(\tau)\,\norm{\errsc(\tau)}
\le \amp\,\norm{\errsc(\tau)}.
\label{eq:thm2-step2}
\end{align}

Here $\amp = \sup_{\tau \in [0, T_{d}]} \amp(\tau)$ is precisely the Score Shocks amplification exponent defined in \citet{sarkar2026scoreshocks}. This is the \textbf{origin} of the $\exp(\amp T_{d})$ factor --- the Lipschitz constant of the drift field $b$ equals $\amp$, so errors propagate at exponential rate $\amp$.

\medskip
\noindent
\textbf{Step~3 --- Controlling $(\mathrm{II})$ (score model error).}
From the definitions of $b, \hat b$ and Assumption~(A2):
\begin{align}
\norm{\hat b(\uhat, \tau) - b(\uhat, \tau)}
&= \abs{\sigtau\,\dot\sigma(\tau)} \cdot \norm{\sth(\uhat, \tau) - \score(\uhat, \tau)} \notag \\
&\le C_{\sigma}\,\norm{\sth(\uhat, \tau) - \score(\uhat, \tau)}.
\label{eq:thm2-step3}
\end{align}

\medskip
\noindent
\textbf{Step~4 --- Expectation inequality.}
Take the Euclidean norm and expectation of~\eqref{eq:thm2-step1-decomp}, and define
\begin{equation}
\delta(\tau) \coloneqq \E\,\norm{\errsc(\tau)}.
\label{eq:thm2-step4-delta}
\end{equation}

By Jensen's inequality $\E\norm{\frac{d}{d\tau}\errsc} \ge \abs{\frac{d}{d\tau}\E\norm{\errsc}}$ (strictly valid under absolute continuity; handled via the chain rule; see Remark~1), combining~\eqref{eq:thm2-step2} and~\eqref{eq:thm2-step3}:
\begin{equation}
\delta'(\tau)
\;\le\; \amp\,\delta(\tau)
\;+\; C_{\sigma}\;\E_{\uhat(\tau)}\!\bigl[\norm{\sth(\uhat, \tau) - \score(\uhat, \tau)}\bigr],
\label{eq:thm2-step4-ineq}
\end{equation}
where $\delta'(\tau) = d\delta/d\tau$ along the reverse time direction ($\tau \searrow 0$).

\begin{remark}
\textbf{(Direction of Jensen's inequality).} Strictly, $\frac{d}{d\tau}\E\norm{\errsc} = \E\bigl[\frac{\errsc}{\norm{\errsc}} \cdot \frac{d \errsc}{d\tau}\bigr] \le \E\norm{\frac{d \errsc}{d\tau}}$, so the inequality direction is correct. On the null set where $\norm{\errsc(\tau)} = 0$, the inequality degenerates to $0 \le \E\norm{\frac{d \errsc}{d\tau}}$, which does not affect the integral bound.
\end{remark}

\medskip
\noindent
\textbf{Step~5 --- Repairing the measure mismatch (core technical step).}

\noindent
\textbf{Problem:} The second expectation in~\eqref{eq:thm2-step4-ineq} is taken under the distribution of $\uhat(\tau)$ (the learned trajectory), whereas the training error Assumption~(A3) is controlled under $\ptau$ (the distribution of the true score trajectory $\unat(\tau)$). The two \textbf{measures differ}.

\noindent
\textbf{Repair via triangle decomposition.} Define
\begin{equation}
g(u, \tau) \coloneqq \norm{\sth(u, \tau) - \score(u, \tau)}.
\label{eq:thm2-step5-g}
\end{equation}
Then
\begin{align}
\E_{\uhat(\tau)}[g(\uhat, \tau)]
&= \E_{\unat(\tau)}[g(\unat, \tau)]
   \;+\;
   \Bigl(\E_{\uhat(\tau)}[g(\uhat, \tau)] - \E_{\unat(\tau)}[g(\unat, \tau)]\Bigr) \notag \\[4pt]
&\le \E_{\unat(\tau)}[g(\unat, \tau)]
   \;+\;
   \abs{\,\E_{\uhat(\tau)}[g(\uhat, \tau)] - \E_{\unat(\tau)}[g(\unat, \tau)]\,}.
\label{eq:thm2-step5-decomp}
\end{align}

By Assumption~(A4) (distributional stability), applied to the function $g(\cdot, \tau)$:
\begin{equation}
\abs{\,\E_{\uhat(\tau)}[g] - \E_{\unat(\tau)}[g]\,}
\;\le\;
C_{\mathrm{stab}} \cdot \Wass{1}\!\bigl(\Law(\uhat(\tau)),\, \Law(\unat(\tau))\bigr) \cdot \Lip(g).
\label{eq:thm2-step5-A4}
\end{equation}

By construction of the synchronous coupling,
\begin{equation}
\Wass{1}\!\bigl(\Law(\uhat(\tau)), \Law(\unat(\tau))\bigr) \le \E\norm{\uhat(\tau) - \unat(\tau)} = \delta(\tau).
\label{eq:thm2-step5-W1}
\end{equation}

By Assumptions~(A5) and~(A1), the Lipschitz constant of $g(\cdot, \tau) = \norm{\sth(\cdot, \tau) - \score(\cdot, \tau)}$ satisfies
\begin{equation}
\Lip(g(\cdot, \tau)) \;\le\; L_{s_{\theta}} + L_{s}(\tau) \;\le\; L_{s_{\theta}} + \sup_{\tau} L_{s}(\tau) \;\eqqcolon\; L_{g} < \infty,
\label{eq:thm2-step5-Lg}
\end{equation}
where finiteness is guaranteed by early-stopping $\tau_{\mathrm{end}} > 0$ (see the remark following Assumption~A1).

Combining these:
\begin{equation}
\E_{\uhat(\tau)}[g(\uhat, \tau)]
\;\le\;
\E_{\unat(\tau)}[g(\unat, \tau)]
\;+\;
C_{\mathrm{stab}}\,L_{g}\;\delta(\tau).
\label{eq:thm2-step5-final}
\end{equation}

Define $C' \coloneqq C_{\mathrm{stab}}\,L_{g}$.

\medskip
\noindent
\textbf{Step~6 --- Assembling the differential inequality.}
Substituting~\eqref{eq:thm2-step5-final} into~\eqref{eq:thm2-step4-ineq}:
\begin{equation}
\boxed{\;
\delta'(\tau)
\;\le\; (\amp + C')\,\delta(\tau)
\;+\; C_{\sigma}\;\E_{\unat(\tau)}\!\bigl[\norm{\sth(\unat, \tau) - \score(\unat, \tau)}\bigr]
\;}.
\label{eq:thm2-step6}
\end{equation}

This is a standard first-order linear differential inequality (Gronwall type), where the distribution of $\unat(\tau)$ is precisely $\ptau$, matching the measure in the training error Assumption~(A3) after Step~5's repair.

\medskip
\noindent
\textbf{Step~7 --- Score error integral bound (Cauchy--Schwarz).}
Define
\begin{equation}
\beta(\tau) \coloneqq \E_{\unat(\tau)}\!\bigl[\norm{\sth(\unat, \tau) - \score(\unat, \tau)}\bigr].
\label{eq:thm2-step7-beta}
\end{equation}

By Cauchy--Schwarz:
\begin{align}
\int_{0}^{T_{d}} \beta(\tau)\,d\tau
&\le \sqrt{T_{d}}\;\left(\int_{0}^{T_{d}} \E_{\unat(\tau)}\!\bigl[\norm{\sth - \score}^{2}\bigr]\,d\tau\right)^{\!1/2} \notag \\
&\le \sqrt{T_{d}}\;\varepsilon,
\label{eq:thm2-step7-cs}
\end{align}
where the last step uses Assumption~(A3). The key point: the expectation in~(A3) is taken under the \textbf{true trajectory distribution} $\ptau = \Law(\unat(\tau))$, which matches the measure in~\eqref{eq:thm2-step6} after the Step~5 repair.

\medskip
\noindent
\textbf{Step~8 --- Gronwall reverse-time integration.}
Apply the Gronwall argument to the differential inequality~\eqref{eq:thm2-step6} over reverse time $[T_{d}, 0]$. Set $\bar\amp \coloneqq \amp + C'$. Rewrite~\eqref{eq:thm2-step6} as
\begin{equation}
\frac{d}{d\tau}\!\bigl[\delta(\tau)\,e^{-\bar\amp \tau}\bigr]
\;\le\; C_{\sigma}\,\beta(\tau)\,e^{-\bar\amp \tau},
\label{eq:thm2-step8-rewrite}
\end{equation}
where the derivative is along the reverse direction ($\tau$ decreasing). Integrating forward over $[0, T_{d}]$ (equivalently: writing the reverse Gronwall in forward form):
\begin{equation}
\delta(0)\,e^{-\bar\amp \cdot 0} - \delta(T_{d})\,e^{-\bar\amp T_{d}}
\;\le\; C_{\sigma} \int_{0}^{T_{d}} \beta(\tau)\,e^{-\bar\amp \tau}\,d\tau.
\label{eq:thm2-step8-int}
\end{equation}

Since the initial conditions coincide at the noise endpoint ($\uhat(T_{d}) \sim p_{T_{d}}$, $\unat(T_{d}) \sim p_{T_{d}}$), we have $\delta(T_{d}) = 0$. Hence
\begin{align}
\delta(0)
&\le C_{\sigma}\,e^{\bar\amp T_{d}} \int_{0}^{T_{d}} \beta(\tau)\,e^{-\bar\amp \tau}\,d\tau \notag \\
&\le C_{\sigma}\,e^{\bar\amp T_{d}} \int_{0}^{T_{d}} \beta(\tau)\,d\tau.
\label{eq:thm2-step8-delta0}
\end{align}

Substituting~\eqref{eq:thm2-step7-cs}:
\begin{equation}
\delta(0) \;\le\; C_{\sigma} \sqrt{T_{d}}\;\varepsilon\;e^{(\amp + C')T_{d}}.
\label{eq:thm2-step8-final}
\end{equation}

\medskip
\noindent
\textbf{Step~9 --- $\Wass{1}$ upper bound and absorption of constants.}
By the synchronous coupling construction, the $\Wass{1}$ distance is bounded above by the $L^{1}$ expectation under the coupling:
\begin{equation}
\Wass{1}\!\bigl(\Law(\uhat(0)),\, \rhotrue\bigr)
\;\le\; \E\norm{\uhat(0) - \unat(0)}
\;=\; \delta(0).
\label{eq:thm2-step9-W1}
\end{equation}

Note: strictly, $\Wass{1}$ above couples $\Law(\uhat(0))$ vs.\ $\Law(\unat(0))$, and $\Law(\unat(0)) = \rhotrue$ (since the true reverse ODE exactly recovers the target distribution). Therefore $\Wass{1}(\Law(\uhat(0)), \rhotrue) \le \delta(0)$ holds.

Substituting~\eqref{eq:thm2-step8-final} and setting $C \coloneqq C_{\sigma} \sqrt{T_{d}}\,e^{C' T_{d}}$:
\begin{equation}
\Wass{1}\!\bigl(\Law(\uhat(0)),\, \rhotrue\bigr)
\;\le\; C\,\varepsilon\,e^{\amp T_{d}}.
\label{eq:thm2-step9-final}
\end{equation}

Observe that $C' = C_{\mathrm{stab}}\,L_{g}$, where $L_{g}$ depends on the schedule and network structure but is independent of $\varepsilon$. In the asymptotic regime $\varepsilon \ll 1$, $e^{C'T_{d}} = \mathcal{O}(1)$ and can be absorbed into the constant $C$. The final form is as stated in~\eqref{eq:thm2-main}.
\end{proof}

\medskip
\noindent
\textbf{PDE Connection --- Stable transfer from numerical solution to true PDE solution.}
The above proof is carried out entirely within the diffusion model domain (state space $\R^{d}$) and yields a $\Wass{1}$ bound for $\Law(\uhat(0))$ approximating $\rhotrue = \Law(\physsol)$. The following ensures that this bound \textbf{stably transfers} to the true PDE solution when the training data originates from a numerical solver, without amplification during physical time propagation.

\begin{proposition}[Kruzhkov $L^{1}$-contractivity]
Let $\mathbf{u}, \mathbf{v}$ be two Kruzhkov entropy solutions of the same scalar conservation law $\dt \mathbf{u} + \dx \flux(\mathbf{u}) = 0$ with initial data $\mathbf{u}_{0}, \mathbf{v}_{0}$ respectively. Then at any physical time $t > 0$,
\begin{equation}
\norm{\mathbf{u}(\cdot, t) - \mathbf{v}(\cdot, t)}_{L^{1}(\Omega)}
\;\le\;
\norm{\mathbf{u}_{0} - \mathbf{v}_{0}}_{L^{1}(\Omega)}.
\label{eq:thm2-pde-L1-contraction}
\end{equation}
(See \citet{kruzhkov1970first}, Theorem~1, or \citet{dafermos2016hyperbolic}, \S\,6.3.)
\end{proposition}

\begin{corollary}[Numerical error does not amplify]
If the training data is generated by a numerical solution $\mathbf{u}_{h}$ satisfying $\norm{\mathbf{u}_{h} - \physsol}_{L^{1}} \le \delta_{\mathrm{num}}$, then under the discrete correspondence $\norm{u} \approx \sqrt{\Delta x}\,\norm{\mathbf{u}}_{L^{2}}$,
\begin{equation}
\norm{\uhat(0) - u_{h}}
\;\le\;
\underbrace{\norm{\uhat(0) - u^{\star}}}_{\text{Theorem~2 bound}\; \mathcal{O}(\varepsilon e^{\amp T_{d}})}
\;+\;
\underbrace{\norm{u^{\star} - u_{h}}}_{\text{numerical error}\; \mathcal{O}(\delta_{\mathrm{num}})}.
\label{eq:thm2-pde-transfer}
\end{equation}
By~\eqref{eq:thm2-pde-L1-contraction}, $\delta_{\mathrm{num}}$ does not grow with physical time $T_{\mathrm{phys}}$ --- this is a core property distinguishing conservation law entropy solutions from general PDEs (e.g., Navier--Stokes $L^{2}$ errors may propagate in time, whereas conservation law $L^{1}$ errors are bounded by the initial data).

\textit{Practical implication:} In experiments, high-resolution Godunov/WENO solvers are used to generate training reference solutions; $\delta_{\mathrm{num}}$ is much smaller than the score-matching error $\varepsilon$ (typical grid resolution $N_{x} \ge 2048$ yields $\delta_{\mathrm{num}} \sim 10^{-4}$), so the Theorem~2 bound is the dominant error term.
\end{corollary}

\medskip
\noindent
\textbf{Remarks on the form and limitations of the bound.}

\begin{remark}[Origin and tightness of $\exp(\amp)$]
\label{rem:exp-origin}
The exponential factor $\exp(\amp T_{d})$ in~\eqref{eq:thm2-main} originates from the Lipschitz constant $\amp = \sup_{\tau} \abs{\sigma \dot\sigma}\,L_{s}(\tau)$ of the drift field $b$. When the target distribution contains shocks (i.e., $\rhotrue$ has Dirac point masses), $\nabla_{u} \score$ blows up as $\mathcal{O}(1/\tau)$ in the interfacial layer, leading to $L_{s}(\tau) \sim 1/\tau$ and hence $\amp$ diverges (or becomes large, depending on the early-stopping $\tau_{\mathrm{end}}$). This exponential amplification is the \textbf{inherent deficiency} of standard score-based diffusion when handling discontinuous solutions.

Theorem~3 reduces this amplification factor to $\mathcal{O}(1)$ via BV-aware parameterization, constituting the core contribution of this paper.
\end{remark}

\begin{remark}[Relation to Score Shocks \citep{sarkar2026scoreshocks}]
\label{rem:score-shocks-relation}
The bound structure of this theorem is consistent with Score Shocks Theorem~6.3 ($\Wass{1} \lesssim \varepsilon \exp(\amp T)$). Our independent derivation supplements the following technical points not detailed in Score Shocks:
\begin{itemize}
\item Formal treatment of the measure mismatch (Step~5, Assumption~A4);
\item Explicit tracing of the Lipschitz constant (Assumption~A1);
\item Bridging to Kruzhkov $L^{1}$-contractivity, enabling the bound to transfer from the diffusion model domain to the PDE domain.
\end{itemize}
\end{remark}

\begin{remark}[The constant $C'$ and the network Lipschitz bound]
\label{rem:C-prime}
The auxiliary constant $C' = C_{\mathrm{stab}}\,L_{g}$ that appears in Steps~5--6 accumulates as $e^{C' T_{d}}$ in the exponent. At practical scales ($T_{d} \sim 1$, $L_{g}$ controlled by spectral normalization or weight decay), $e^{C' T_{d}} = \mathcal{O}(1)$, so this constant does not affect the asymptotic rate but only the absolute constant $C$. Formally, one may require the network to satisfy $\norm{\sth}_{\mathrm{Lip}} \le L_{\mathrm{net}}$ as part of training regularization.
\end{remark}


% ----- A1.3 · Theorem 3 · Improved rate (main result) ⭐ -----
\subsection{Proof of Theorem~\ref{thm:improved-rate} (Improved rate)}
\label{appx:proofs:thm3}
%==============================================================================
% proofs/thm3_proof.tex — Theorem 3 full proof (Improved Rate, NeurIPS 2026)
%
% Corresponds to main body \ref{thm:improved-rate}. Six-stage complete proof,
% with assumption system A0--A6.
% Macros from macros/notation.tex.
%==============================================================================

% ============================================================================
% Theorem Statement
% ============================================================================
\begin{theorem}\label{appx:thm3}
Under Assumptions~\ref{ass:A0}--\ref{ass:A6}, consider the BV-aware parameterization~\eqref{eq:bv-aware} and the loss $\mathcal{L}_3$ (including the TV regularizer).
Let the score matching error satisfy A3, with $\varepsilon \in (0,1)$.
Then the generated distribution $\muth = \Law(\uth)$ and the Dirac measure $\delta_{\physsol}$ at the Kruzhkov entropy solution $\physsol$ satisfy
\begin{equation}\label{eq:thm3-bound}
  \Wass{1}(\muth, \delta_{\physsol}) \leq C_{\mathrm{final}} \, \varepsilon^{1/2},
\end{equation}
where $C_{\mathrm{final}}$ depends only on $M, T_d, T_{\mathrm{phys}}, \kappa_0, L_{\mathrm{sm}}, \physvis/\varepsilon$ (all independent of $\varepsilon$).
In particular, the bound \emph{does not} contain the exponential amplification factor $\exp(\amp T)$ (cf.\ Theorem~\ref{thm:stability}).
\end{theorem}

% ============================================================================
% Assumption System
% ============================================================================
\subsection{Assumption System (A0--A6)}\label{sec:assumptions}

The following assumptions hold throughout the proof. The physical meaning and verifiability of each are stated alongside.

\begin{description}
\item[Assumption A0 (BV boundedness)]\label{ass:A0}
The TV regularizer $\lambda_{\mathrm{BV}} \, \TV(\Dth)$ in the loss $\mathcal{L}_3$ together with the hard-coded $\tanh$ structure of parameterization~\eqref{eq:bv-aware} jointly guarantee that the generated sample $\uth$ almost surely belongs to the bounded TV ball
\begin{equation}
  \mathcal{K}_M \coloneqq \bigl\{ \mathbf{u} \in L^1(\Omega) \cap L^\infty(\Omega) : \TV(\mathbf{u}) \leq M \bigr\},
\end{equation}
and the true entropy solution satisfies $\TV(\physsol) \leq M$, with $M > 0$ a constant independent of $\varepsilon$.
\emph{Verifiability:} $M$ can be estimated from the TV bound of the initial data and the TV-non-increasing property of the BV semigroup; $\lambda_{\mathrm{BV}}$ is a training hyperparameter.

\item[Assumption A1 (Locally bimodal score structure)]\label{ass:A1}
In a neighborhood of the shock manifold $\Shockscore(\tau)$, the true score field admits the asymptotic form
\begin{equation}
  \score(u, \tau) = s^{\mathrm{sm}}(u, \tau) + \frac{\Delta s_{\mathrm{jump}}(u)}{2} \tanh\!\Bigl(\frac{A(u,\tau) - B(u,\tau)}{2\tau}\Bigr) + \mathcal{O}(\tau),
\end{equation}
where $A, B$ are the Cole--Hopf potential functions of the two dominant modes, $\Delta s_{\mathrm{jump}} = \nabla_u B - \nabla_u A$ is the jump vector, and $s^{\mathrm{sm}}$ is uniformly Lipschitz on the whole space.
This structure follows as a local result of the Laplace asymptotic method, assuming sufficiently small diffusion time and well-separated Gaussian clusters; superposition of finitely many such local structures is allowed.
\emph{Verifiability:} Verifiable for discrete distribution mixture models; treated as a reasonable engineering approximation for general continuous distributions (see engineering approximation EA1).

\item[Assumption A2 (Network parameterization constraint)]\label{ass:A2}
In the parameterization~\eqref{eq:bv-aware}:
\begin{itemize}
  \item $\jumpamp(u,\tau) \geq \kappa_0 > 0$ holds strictly for all $(u,\tau)$ (enforceable via Softplus activation, i.e., $\jumpamp = \kappa_0 + \mathrm{Softplus}(\tilde{\kappa}_\theta)$);
  \item The maps $\nabla_u^2 \phisbg$, $\nabla_u^2 \Dth$, $\nabla_u \jumpamp$ are uniformly bounded on $(u,\tau) \in \R^d \times [\tau_{\mathrm{end}}, T_d]$, with bounds independent of $\tau$.
\end{itemize}
\emph{Verifiability:} Enforceable at training time via network architecture design and weight norm constraints.

\item[Assumption A3 (Score matching error)]\label{ass:A3}
There exists $\varepsilon > 0$ such that in the grid-independent mean-squared error sense:
\begin{equation}
  \E_{u,\tau}\!\Bigl[ \frac{|\Omega|}{d} \sum_{i=1}^d \errsc_i(u,\tau)^2 \Bigr] \leq \varepsilon^2,
\end{equation}
where $\Delta x = |\Omega|/d$; equivalently $\E_{u,\tau}\bigl[\Delta x \norm{\errsc(u,\tau)}^2\bigr] \leq \varepsilon^2$, with expectation taken over the training distribution $p(u,\tau)$ ($u \sim p(\cdot,\tau)$, $\tau \sim \mathcal{U}[\tau_{\mathrm{end}}, T_d]$).
This setting makes $\varepsilon^2$ a physical approximation of the continuous $\int_\Omega \errsc(x,\tau)^2 \, dx$.

\item[Assumption A4 (Prior matching)]\label{ass:A4}
The squared distance bound of the prior approximation at the initial time $\tau = T_d$ of the reverse ODE is controlled by the same scaling:
\begin{equation}
  \E\bigl[\Delta x \norm{\uhat(T_d) - \unat(T_d)}^2 \bigr] \leq \tilde{C}^2 \varepsilon^2,
\end{equation}
where $\tilde{C} > 0$ is a constant independent of $\varepsilon$ and $d$.
\emph{Typically satisfied when:} $p_{T_d}$ is standard Gaussian and $T_d$ is large enough that the signal-to-noise ratio is very low.

\item[Assumption A5 (Early stopping and viscosity matching balance, engineering approximation EA2)]\label{ass:A5}
The reverse ODE is early-stopped at $\tau_{\mathrm{end}} > 0$, the effective terminal viscosity is defined by $\sigtau^2(\tau_{\mathrm{end}}) = 2\physvis \tau_{\mathrm{end}}$, and the following holds:
\begin{equation}
  \physvis = \Theta(\varepsilon), \quad \text{i.e.,} \quad c\varepsilon \leq \physvis \leq C\varepsilon,
\end{equation}
with A3 still valid over the full interval $[\tau_{\mathrm{end}}, T_d]$.
\emph{Engineering approximation note:} The achievability of $\physvis = \Theta(\varepsilon)$ depends on fine design of the noise schedule and prior estimation of the training error $\varepsilon$; this is an engineering approximation (EA2), whose scope of influence is discussed in Appendix~C.

\item[Assumption A6 (Kuznetsov-type error bound)]\label{ass:A6}
Under A0 and A5, the $L^1$ error between the viscous PDE solution $\physsolvis$ and the Kruzhkov entropy solution $\physsol$ satisfies
\begin{equation}
  \norm{\physsolvis - \physsol}_{L^1(\Omega)} \leq C_K \sqrt{\physvis},
\end{equation}
where $C_K$ depends on the initial TV bound $M$ and physical time $T_{\mathrm{phys}}$, but is independent of $\varepsilon$.
\emph{Explanation:} This bound follows directly from the classical theorem of Kuznetsov (1976) under A0, with $C_K = C \sqrt{T_{\mathrm{phys}}} \cdot \TV(\initdat) \leq C \sqrt{T_{\mathrm{phys}}} \cdot M$.
\end{description}

% ============================================================================
% Proof
% ============================================================================
\begin{proof}

The proof proceeds in four stages:
\begin{enumerate}
  \item[Stage~1.] Using A1 and A2, construct a function-space decomposition and analyze the well-posedness of the BV-aware parameterization.
  \item[Stage~2.] Based on one-sided Lipschitz analysis and a modified Gronwall inequality, using A2, A3, A4 to establish the trajectory error bound.
  \item[Stage~3.] Introduce the physical PDE framework and Kruzhkov theory, stating the BV compactness required for Stage~4.
  \item[Stage~4.] Using A0, A5, A6 to assemble the Wasserstein error, obtaining $\Wass{1} \leq C_{\mathrm{final}} \varepsilon^{1/2}$.
\end{enumerate}

% ----------------------------------------------------------------------------
% Preliminaries: Notation Summary
% ----------------------------------------------------------------------------
\subsection*{Notation (unified for the entire proof)}\label{sec:notation-proof}

\paragraph{State space (diffusion model domain)}
\begin{description}
  \item[$u \in \R^d$] State vector ($d = N_x$ is the number of spatial grid points)
  \item[{$\tau \in [0, T_d]$}] Diffusion time (reverse: $T_d \to 0$, i.e., noise $\to$ data)
  \item[$p(u, \tau)$] Smoothed state density
  \item[$\score(u, \tau) = \nabla_u \log p(u, \tau)$] True score field
  \item[$\sth(u, \tau)$] Network-approximated score field
  \item[$\errsc(u, \tau) = \sth(u, \tau) - \score(u, \tau)$] Score residual field
  \item[$\Shockscore(\tau) \subset \R^d$] Shock manifold
  \item[$\dShock(u, \tau)$] Signed distance function to $\Shockscore(\tau)$
\end{description}

\paragraph{Physical space (PDE domain)}
\begin{description}
  \item[$x \in \Omega \subset \R$] Physical space coordinate
  \item[{$t \in [0, T]$}] Physical time (evolved by the PDE from initial data, orthogonal to $\tau$)
  \item[$\physsol(x)$] Kruzhkov entropy solution (at fixed $t = T_{\mathrm{phys}}$)
  \item[$\physsolvis(x)$] Viscous PDE approximation with viscosity $\physvis$ (at $t = T_{\mathrm{phys}}$)
  \item[$\dx$] Physical space partial derivative
  \item[$\norm{\cdot}_{L^1(\Omega)}, \norm{\cdot}_{L^2(\Omega)}$] $L^p$ norms on physical space
\end{description}

\paragraph{Trajectories and generation}
\begin{description}
  \item[$\uhat(\tau)$] ODE trajectory driven by the learned score ($\R^d$-valued)
  \item[$\unat(\tau)$] ODE trajectory driven by the true score ($\R^d$-valued)
  \item[$\uth \coloneqq \uhat(0)$] Final generated sample (at $\tau = 0$)
  \item[$\muth \coloneqq \Law(\uth)$] Distribution of the generated samples
  \item[$\delta_{\physsol}$] Dirac measure at the target Kruzhkov entropy solution
\end{description}

\paragraph{Norms and distances}
\begin{description}
  \item[$\norm{v}$] Euclidean norm on state space $\R^d$
  \item[$\langle \cdot, \cdot \rangle$] $\R^d$ inner product
  \item[$\TV(f) = \int_\Omega |\dx f| \, dx$] Total variation seminorm
  \item[$\Wass{1}(\mu, \nu)$] 1-Wasserstein distance (Kantorovich--Rubinstein dual)
\end{description}

\paragraph{Time mapping}
The physical meaning of diffusion time $\tau$ decreasing from $T_d$ to $0$:
\begin{itemize}
  \item $\tau = T_d$: pure noise (prior, e.g., standard Gaussian)
  \item $\tau = 0$: data sample (corresponding to the PDE solution at physical time $t = T_{\mathrm{phys}}$)
  \item \textbf{Viscosity matching}: $\sigtau^2(\tau) = 2\physvis \tau$, so the evolution of $\tau$ is equivalent to a physical process with viscosity $\physvis$; the drift coefficient of the reverse ODE is $\physvis$
  \item In Stage~3, the target $\physsol$ is the PDE entropy solution at a \emph{fixed physical time} $t = T_{\mathrm{phys}}$ ($T_{\mathrm{phys}}$ is regarded as a constant and is not written explicitly in the sequel)
  \item Diffusion time $\tau$ and physical time $t$ are \emph{orthogonal}: $\tau$ describes the noising--denoising intensity, $t$ describes the PDE evolution
\end{itemize}

% ============================================================================
% Stage 1: Function Space and Error Localization
% ============================================================================
\subsection{Stage 1: Function Space and Error Localization}\label{sec:stage1}

\subsubsection{1.1 Setup}

Let the dimensionless diffusion time of the target distribution be $\tau \in (\tau_{\mathrm{end}}, T_d]$. Denote the true smoothed data density by
\begin{equation}
  p(u, \tau) = (\rho \ast \Gtau)(u), \qquad \Gtau(u) = (4\pi\tau)^{-d/2} \exp\!\bigl(-\norm{u}^2/4\tau\bigr),
\end{equation}
where $\rho(u)$ is the target data distribution. The true score field:
\begin{equation}
  \score(u, \tau) = \nabla_u \log p(u, \tau).
\end{equation}

Define the shock manifold $\Shockscore(\tau) \subset \R^d$ as the equipotential surface in state space that connects different modes. For any $u \in \R^d$, define its signed distance to $\Shockscore(\tau)$:
\begin{equation}
  \dShock(u, \tau) = \inf_{v \in \Shockscore(\tau)} \norm{u - v} \cdot \mathrm{sgn}(\text{mode label}).
\end{equation}

The shock interfacial layer neighborhood and the outer smooth region:
\begin{equation}
  \Omega_{\mathrm{sh}}(\tau) = \bigl\{ u \in \R^d \;\big|\; |\dShock(u, \tau)| \leq \mathcal{O}(\tau) \bigr\}, \qquad
  \Omega_{\mathrm{out}}(\tau) = \R^d \setminus \Omega_{\mathrm{sh}}(\tau).
\end{equation}

\subsubsection{1.2 Lemma 1.1: Asymptotic Expansion and Singularity Decomposition of the True Score Field}

\begin{lemma}\label{lem:1.1}
Under Assumption~A1, the true score field $\score(u, \tau)$ can be decomposed on the whole space $\R^d$ as
\begin{equation}
  \score(u, \tau) = s^{\mathrm{sm}}(u, \tau) + s^{\mathrm{sing}}(u, \tau) + \mathcal{O}(\tau),
\end{equation}
where:
\begin{enumerate}
  \item $s^{\mathrm{sm}}(u, \tau)$ is uniformly Lipschitz continuous on the whole space: $\norm{\nabla_u s^{\mathrm{sm}}}_{\mathrm{op}} \leq L_{\mathrm{sm}} = \mathcal{O}(1)$.
  \item $s^{\mathrm{sing}}(u, \tau)$ is non-negligible only inside $\Omega_{\mathrm{sh}}(\tau)$, and has the exact form
  \begin{equation}
    s^{\mathrm{sing}}(u, \tau) = \frac{\Delta s_{\mathrm{jump}}(u)}{2} \tanh\!\Bigl(\frac{A(u,\tau) - B(u,\tau)}{2\tau}\Bigr),
  \end{equation}
  where $A, B$ are the potential functions of the dominant modes and $\Delta s_{\mathrm{jump}} = \nabla_u B - \nabla_u A$ is the jump intensity vector.
\end{enumerate}
\end{lemma}

\textit{Proof sketch (Laplace method):} By A1, inside $\Omega_{\mathrm{sh}}(\tau)$, the integral representation of $p(u,\tau)$ is dominated by the two nearest modes:
\begin{equation}
  p(u, \tau) \sim C_A(u) \exp\!\Bigl(-\frac{A(u,\tau)}{\tau}\Bigr) + C_B(u) \exp\!\Bigl(-\frac{B(u,\tau)}{\tau}\Bigr).
\end{equation}
Taking the log-gradient and using the identity $\nabla_u \log(e^{-A/\tau} + e^{-B/\tau})$ extracts the $\tanh$ part; the slowly varying coefficients $C_A, C_B$ and their gradients are absorbed into $s^{\mathrm{sm}}$.
In $\Omega_{\mathrm{out}}(\tau)$, one mode is overwhelmingly dominant, $\tanh(\cdot) \to \pm 1$, $\nabla_u s^{\mathrm{sing}}$ tends to $0$, and the score field degenerates to the smooth $\nabla_u A$ or $\nabla_u B$.
The error term $\mathcal{O}(\tau)$ originates from the next-order remainder of the Laplace method; the uniform estimate is guaranteed by the mode-separation condition of A1.\hfill$\blacksquare$

\begin{remark}[Nature of the singularity]
Under standard parameterization, if the network directly fits $\score(u,\tau)$, the Jacobian $\nabla_u s^{\mathrm{sing}}$ of $s^{\mathrm{sing}}$ contains a singular component of $\mathcal{O}(1/\tau)$ inside $\Omega_{\mathrm{sh}}$---this is precisely the origin of the global Lipschitz constant $L(\tau) \propto 1/\tau$ in classical theory, which causes the Gronwall integral to diverge and produces the exponential amplification factor $\exp(\amp T_d)$.
The core idea of parameterization~\eqref{eq:bv-aware} is to \emph{explicitly embed} this singularity into the network structure, so that the residual field $\errsc(u,\tau)$ no longer inherits the $1/\tau$ singularity.
\end{remark}

\subsubsection{1.3 Lemma 1.2: Well-Posedness of the BV-Aware Parameterization}

The parameterization~\eqref{eq:bv-aware} defines the network score field:
\begin{equation}
  \sth(u, \tau) = \underbrace{\nabla_u \phisbg(u, \tau)}_{\text{smooth background term}} + \underbrace{\frac{\jumpamp(u,\tau)}{2} \tanh\!\Bigl(\frac{\Dth(u, \tau)}{2\tau}\Bigr) \nabla_u \Dth(u, \tau)}_{\text{interfacial singular term}}.
\end{equation}

\begin{lemma}\label{lem:1.2}
(Spatial alignment and network well-posedness) Under Assumption~A2, the map $\Theta: (u,\tau) \mapsto (\phisbg, \jumpamp, \Dth)$ has the following properties:
\begin{enumerate}
  \item Target approximation: $\nabla_u \phisbg \approx s^{\mathrm{sm}}$, $\jumpamp \nabla_u \Dth \approx \Delta s_{\mathrm{jump}}$, $\Dth \approx A-B$.
  \item Bounded Jacobian: The Jacobian norms $\norm{\nabla_u^2 \phisbg}_{\mathrm{op}}$, $\norm{\nabla_u \jumpamp}$, $\norm{\nabla_u^2 \Dth}_{\mathrm{op}}$ are all independent of $\tau$, of order $\mathcal{O}(1)$.
  \item The network learns only the smooth signed distance field $\Dth$. The singular factor $1/\tau$ is written as an explicit physical prior \emph{outside} the network structure, appearing only inside $\tanh$ and in the singular Jacobian matrix $J_{\mathrm{sh}}$ after differentiation; the network involves no $1/\tau$ amplification from automatic differentiation, completely avoiding gradient explosion.
\end{enumerate}
\end{lemma}

\textit{Proof:} By A2, $\nabla_u^2 \phisbg$, $\nabla_u^2 \Dth$, $\nabla_u \jumpamp$ are all uniformly bounded. Note that the derivative of $\tanh$ with respect to its argument is $\operatorname{sech}^2(\cdot) \leq 1$, so when computing $\nabla_u$ of $\sth$, the dependence of $\tau$ in the denominator does not propagate divergently outward; it passes into the Jacobian and forms a strictly negative semi-definite singular part (proved in detail below).
Since $A$ and $B$ are both $\mathcal{O}(1)$ potential functions, the network-learned $\Dth(u, \tau)$ has gradient $\nabla_u \Dth = \mathcal{O}(1)$, introducing no additional singular scale.\hfill$\blacksquare$

\subsubsection{1.4 Boundedness of the Residual Field (Heuristic Discussion)}

\begin{remark}[Heuristic significance of Corollary 1.4; not a rigorous claim]
The following is a heuristic argument for intuition only and is \emph{not relied upon in subsequent steps of the proof}.

Define the score residual field $\errsc(u,\tau) = \sth(u,\tau) - \score(u,\tau)$, decomposed into $\errsc = e_{\mathrm{sm}} + e_{\mathrm{sh}}$, where $e_{\mathrm{sh}}$ is the shock residual arising from interface location/intensity deviation.

In the standard architecture, $e_{\mathrm{standard}}$ cannot fit the steep gradient of $\tanh(\cdot/\tau)$ and itself has a gradient component of $-\mathcal{O}(1/\tau)$ inside $\Omega_{\mathrm{sh}}$, causing exponential error amplification.
The BV-aware architecture embeds the $\tanh$ structure, making the error of $e_{\mathrm{sh}}$ manifest as a small shift of the shock center position; this shift is of the same order as the shift magnitude in the $\Wass{1}$ sense ($\Delta \Wass{1} \propto \Delta x$), without involving multiplicative amplification of gradients.

More rigorously, if one can independently prove gradient bounds for $e_{\mathrm{sm}}$ and for the shock shift of $e_{\mathrm{sh}}$, one can establish
$\int_0^{T_d} \norm{\nabla_u \errsc(u,\tau)}_{\mathrm{eff}} \, d\tau \leq C_0 = \mathcal{O}(1)$.
Rigorization of this direction is left to future work (see technical point C1.4-G in Appendix~C).
In the present proof, the Gronwall argument of Stage~2 \emph{does not rely on this bound} and uses only the integral estimate of $\norm{\errsc}$ itself (see Theorem~2.2).
\end{remark}

% ============================================================================
% Stage 2: Modified Gronwall Inequality and Trajectory Error Bound
% ============================================================================
\subsection{Stage 2: Modified Gronwall Inequality and Trajectory Error Bound}\label{sec:stage2}

\subsubsection{2.1 Failure of the Classical Analysis and the Gronwall Trap}

In the reverse generative process, consider two ODE trajectories---the reference trajectory $\unat(\tau)$ driven by the true score $\score$ and the approximate trajectory $\uhat(\tau)$ driven by the network score $\sth$.
Both obey the probability-flow ODE under the viscosity matching strategy $\sigtau^2(\tau) = 2\physvis \tau$ (drift coefficient $\sigma \dot\sigma = \physvis$):
\begin{equation}\label{eq:ode-pair}
  \frac{d \unat}{d\tau} = -\physvis \, \score(\unat, \tau) \eqqcolon V(\unat, \tau), \qquad
  \frac{d \uhat}{d\tau} = -\physvis \, \sth(\uhat, \tau) \eqqcolon V_\theta(\uhat, \tau).
\end{equation}

\textbf{Initial condition (A4).} By A4, the error between the two trajectories at $\tau = T_d$ satisfies
\begin{equation}
  E(T_d) \coloneqq \norm{\uhat(T_d) - \unat(T_d)} \leq \tilde{C}\varepsilon,
\end{equation}
(equal to $0$ if the prior matches exactly; small deviations are handled uniformly by the Gronwall formula below).

Define the trajectory error:
\begin{equation}
  E(\tau) \coloneqq \norm{\uhat(\tau) - \unat(\tau)}.
\end{equation}

Differentiating $E^2$:
\begin{equation}
  \frac{1}{2}\frac{d}{d\tau} E(\tau)^2
  = \bigl\langle \uhat - \unat,\; V_\theta(\uhat, \tau) - V(\unat, \tau) \bigr\rangle.
\end{equation}

Decompose the right-hand side into the self-consistency term (of $V_\theta$) and the error-driven term:
\begin{align}
  \frac{1}{2}\frac{d}{d\tau} E(\tau)^2
  &\leq \underbrace{\bigl\langle \uhat - \unat,\; V_\theta(\uhat,\tau) - V_\theta(\unat,\tau) \bigr\rangle}_{\text{Term~I: drift self-consistency}} \\
  &\quad + \underbrace{\physvis \bigl\langle \uhat - \unat,\; \errsc(\unat,\tau) \bigr\rangle}_{\text{Term~II: score error-driven}}.\notag
\end{align}

Applying Cauchy--Schwarz to Term~II:
$\physvis \langle \uhat - \unat, \errsc(\unat,\tau) \rangle \leq \physvis \norm{\errsc(\unat,\tau)} \cdot E(\tau)$.
Thus:
\begin{equation}
  \frac{1}{2}\frac{d}{d\tau} E(\tau)^2
  \leq \underbrace{\bigl\langle \uhat - \unat,\; V_\theta(\uhat,\tau) - V_\theta(\unat,\tau) \bigr\rangle}_{\text{Term~I}} + \physvis \norm{\errsc(\unat,\tau)} \cdot E(\tau).
\end{equation}

In the classical proof, Term~I is bounded using the global Lipschitz constant $L_\theta(\tau) = \mathcal{O}(1/\tau)$, causing the Gronwall integral to diverge as $\exp(\amp T_d)$.
Below we replace this with a one-sided Lipschitz constant.

\subsubsection{2.2 Lemma 2.1: One-Sided Lipschitz Condition and Negative Semi-Definiteness}

\begin{definition}[One-sided Lipschitz constant]
The one-sided Lipschitz constant $\mu(\tau)$ of $V_\theta$ is the smallest number such that for all $u, v \in \R^d$:
\begin{equation}
  \bigl\langle u - v,\; V_\theta(u,\tau) - V_\theta(v,\tau) \bigr\rangle \leq \mu(\tau) \norm{u - v}^2.
\end{equation}
When $V_\theta$ is differentiable, $\mu(\tau)$ equals the maximum eigenvalue of the symmetric part of its Jacobian $J_\theta(u,\tau) \coloneqq \nabla_u V_\theta(u,\tau)$:
\begin{equation}
  \mu(\tau) = \sup_{u \in \R^d} \lambda_{\max}\!\Bigl(\frac{J_\theta(u,\tau) + J_\theta(u,\tau)^\mathsf{T}}{2}\Bigr).
\end{equation}
\end{definition}

\begin{lemma}\label{lem:2.1}
(Bounded logarithmic norm of the BV-aware architecture) Under Assumption~A2,
\begin{equation}
  \mu(\tau) \leq L_{\mathrm{sm}} = \mathcal{O}(1),
\end{equation}
i.e., the one-sided Lipschitz constant is independent of $\tau$ and has no singularity.
\end{lemma}

\textit{Proof:} Compute the Jacobian of $V_\theta(u,\tau) = -\physvis \sth(u,\tau)$:
\begin{equation}
  J_\theta = \nabla_u V_\theta = J_{\mathrm{sm}} + J_{\mathrm{sh}} + J_{\mathrm{cross}},
\end{equation}
with the following components.

\textbf{(a) Smooth part.} $J_{\mathrm{sm}} = -\physvis \nabla_u^2 \phisbg$. By A2, $\norm{\nabla_u^2 \phisbg}_{\mathrm{op}} = \mathcal{O}(1)$, so $\lambda_{\max}(J_{\mathrm{sm}}) \leq \physvis \cdot \mathcal{O}(1) = \mathcal{O}(1)$ (since $\physvis = \Theta(\varepsilon)$ is small but $\varepsilon \leq 1$, so it does not exceed $\mathcal{O}(1)$).

\textbf{(b) Cross terms.} $J_{\mathrm{cross}}$ contains derivative terms of $\jumpamp$ and $\nabla_u \Dth$. By A2 these are all uniformly bounded; moreover, the dominant terms produced by these derivatives are already controlled by explicit separation and internal derivative control, introducing no $1/\tau$ factor into this term, so $\lambda_{\max}(J_{\mathrm{cross}}) = \mathcal{O}(1)$.

\textbf{(c) Singular part.} Differentiating the $\tanh$ term in the parameterization with respect to $\nabla_u$ yields
\begin{equation}
  J_{\mathrm{sh}}(u,\tau) = -\physvis \cdot \frac{\jumpamp(u,\tau)}{4\tau} \, \operatorname{sech}^2\!\Bigl(\frac{\Dth(u,\tau)}{2\tau}\Bigr) \, \bigl(\nabla_u \Dth\bigr)\bigl(\nabla_u \Dth\bigr)^\mathsf{T}.
\end{equation}
The matrix $(\nabla_u \Dth)(\nabla_u \Dth)^\mathsf{T}$ is a rank-$1$ positive semi-definite matrix (of the form $\mathbf{n}\mathbf{n}^\mathsf{T}$, where $\mathbf{n} = \nabla_u \Dth$ is the shock interface normal vector).
By A2, $\jumpamp \geq \kappa_0 > 0$; together with $\physvis > 0$, $\tau > 0$, $\operatorname{sech}^2(\cdot) > 0$, the leading negative sign makes $J_{\mathrm{sh}}$ \emph{strictly negative semi-definite}.

Hence the maximum eigenvalue of the symmetric part of $J_{\mathrm{sh}}$ satisfies $\lambda_{\max}(J_{\mathrm{sh}}) \leq 0$, contributing nothing positive to the one-sided Lipschitz constant.

\textbf{(d) Physical interpretation.} In the direction normal to the shock surface (the $\mathbf{n}$ direction), $J_{\mathrm{sh}}$ provides a negative eigenvalue of magnitude $-\physvis \kappa_0 \norm{\nabla_u \Dth}^2 / (4\tau)$, corresponding to a \emph{strong attractive force} of the reverse diffusion flow in the shock normal direction. Errors that deviate from the true shock manifold are compressed at this enormous rate rather than amplified---this is the mathematical essence of Burgers characteristic line convergence, and the core guarantee that this architecture avoids exponential divergence.

Combining (a)--(d),
\begin{equation}
  \mu(\tau) \leq \lambda_{\max}\!\bigl(J_{\mathrm{sm}} + J_{\mathrm{cross}}\bigr) + \underbrace{\lambda_{\max}(J_{\mathrm{sh}})}_{\leq 0} \leq L_{\mathrm{sm}} = \mathcal{O}(1).
\end{equation}
\hfill$\blacksquare$

\subsubsection{2.3 Theorem 2.2: Modified Flow Stability with Constant Upper Bound}

\begin{lemma}\label{lem:2.2}
(Trajectory error bound) Under Assumptions~A2, A3, A4, let $\varepsilon_{\mathrm{raw}}$ be the raw (unnormalized) score matching mean-squared error bound, i.e., $\E[\norm{\errsc}^2] \leq \varepsilon_{\mathrm{raw}}^2$.
By the grid-normalized definition of A3, $\E[\Delta x \norm{\errsc}^2] \leq \varepsilon^2$, so $\varepsilon_{\mathrm{raw}}^2 = \varepsilon^2 / \Delta x$. Then
\begin{equation}
  \E_z\bigl[E(0)\bigr] \leq C_1 \varepsilon_{\mathrm{raw}} = C_1 \frac{\varepsilon}{\sqrt{\Delta x}},
\end{equation}
where $C_1$ depends only on $T_d, \physvis, L_{\mathrm{sm}}$. In the norm conversion of Stage~4, $\sqrt{\Delta x}$ will be canceled, yielding a grid-independent final bound.
\end{lemma}

\textit{Proof:}
\textbf{Step 1: Scalar differential inequality.}
Applying Lemma~2.1 to Term~I in the derivative of $\frac{1}{2}E^2$ gives
\begin{equation}
  \frac{1}{2}\frac{d}{d\tau} E(\tau)^2 \leq \mu(\tau) \, E(\tau)^2 + \physvis \, \norm{\errsc(\unat,\tau)} \cdot E(\tau).
\end{equation}
Dividing by $E(\tau)$ (when $E(\tau) > 0$; the inequality holds trivially when $E(\tau) = 0$),
\begin{equation}\label{eq:star}
  \frac{d}{d\tau} E(\tau) \leq \mu(\tau) \, E(\tau) + \physvis \, \norm{\errsc(\unat,\tau)}. \tag{$\star$}
\end{equation}

\textbf{Step 2: Apply Gronwall's inequality.}
From $(\star)$, integrating from $\tau = T_d$ to $\tau = 0$,
\begin{equation}
  E(0) \leq E(T_d) \cdot \exp\!\Bigl(\int_0^{T_d} \mu(r) \, dr\Bigr) + \int_0^{T_d} \exp\!\Bigl(\int_0^\tau \mu(r) \, dr\Bigr) \, \physvis \, \norm{\errsc(\unat,\tau)} \, d\tau.
\end{equation}

By Lemma~2.1, $\mu(r) \leq L_{\mathrm{sm}} = \mathcal{O}(1)$, so
\begin{equation}
  \exp\!\Bigl(\int_0^\tau \mu(r) \, dr\Bigr) \leq e^{L_{\mathrm{sm}} T_d} \eqqcolon C_{\mathrm{flow}} = \mathcal{O}(1).
\end{equation}

Substituting the initial error bound $E(T_d) \leq \tilde{C}\varepsilon$ from A4:
\begin{equation}\label{eq:starstar}
  E(0) \leq \tilde{C} C_{\mathrm{flow}} \, \varepsilon + C_{\mathrm{flow}} \, \physvis \int_0^{T_d} \norm{\errsc(\unat,\tau)} \, d\tau. \tag{$\star\star$}
\end{equation}

\textbf{Step 3: Integral bound of the score matching error (directly from A3).}
Apply Cauchy--Schwarz to the integral $\int_0^{T_d} \norm{\errsc(\unat,\tau)} \, d\tau$:
\begin{equation}
  \int_0^{T_d} \norm{\errsc(\unat,\tau)} \, d\tau \leq \sqrt{T_d} \cdot \Bigl(\int_0^{T_d} \norm{\errsc(\unat,\tau)}^2 \, d\tau\Bigr)^{1/2}.
\end{equation}

Take expectation over $z$ (prior noise) of the right-hand side. By A3, $\E[\Delta x \norm{\errsc}^2] \leq \varepsilon^2$, i.e., the raw (unnormalized) mean-squared error satisfies $\E[\norm{\errsc}^2] \leq \varepsilon^2 / \Delta x = \varepsilon_{\mathrm{raw}}^2$. Hence
\begin{equation}
  \E_z\!\Bigl[\int_0^{T_d} \norm{\errsc(\unat,\tau)}^2 \, d\tau\Bigr] \leq T_d \cdot \varepsilon_{\mathrm{raw}}^2 = T_d \cdot \frac{\varepsilon^2}{\Delta x},
\end{equation}
(this step holds because the marginal distribution of $\unat(\tau)$ is $p(\cdot,\tau)$, so the expectation over $(\unat(\tau),\tau)$ matches $\E_{u,\tau}$ in A3).

Now take expectation of the Cauchy--Schwarz result (using Jensen's inequality):
\begin{equation}
  \E_z\!\Bigl[\int_0^{T_d} \norm{\errsc(\unat,\tau)} \, d\tau\Bigr] \leq \sqrt{T_d} \cdot \sqrt{T_d \, \varepsilon_{\mathrm{raw}}^2} = T_d \, \varepsilon_{\mathrm{raw}} = T_d \frac{\varepsilon}{\sqrt{\Delta x}}.
\end{equation}

\textbf{Step 4: Assembly.}
Substituting Step~3 into $(\star\star)$ and taking expectation:
\begin{equation}
  \E_z[E(0)] \leq \tilde{C} C_{\mathrm{flow}} \, \varepsilon_{\mathrm{raw}} + C_{\mathrm{flow}} \, \physvis \, T_d \, \varepsilon_{\mathrm{raw}} \eqqcolon C_1 \varepsilon_{\mathrm{raw}},
\end{equation}
where $C_1 = C_{\mathrm{flow}}(\tilde{C} + \physvis T_d) = \mathcal{O}(1)$ (by A5, $\physvis = \Theta(\varepsilon) \leq 1$, so $\physvis T_d = \mathcal{O}(1)$).
From $\varepsilon_{\mathrm{raw}} = \varepsilon / \sqrt{\Delta x}$, we obtain $\E_z[E(0)] \leq C_1 \varepsilon / \sqrt{\Delta x}$.
\hfill$\blacksquare$

\textbf{Conclusion.} Under the BV-aware architecture, the trajectory endpoint error is $\E[E(0)] = \mathcal{O}(\varepsilon_{\mathrm{raw}})$, free of the exponential amplification factor.

% ============================================================================
% Stage 3: Kruzhkov Entropy Solution L^1 Contraction Framework
% ============================================================================
\subsection{Stage 3: Kruzhkov Entropy Solution \texorpdfstring{$L^1$}{L1} Contraction Framework}\label{sec:stage3}

\begin{remark}[Convention for Stage~3]
We now introduce physical space $x \in \Omega \subset \R$. The components of the state vector $u \in \R^d$ ($d = N_x$) are $u_i = \mathbf{u}(x_i)$, where $\{x_i\}$ is a uniform grid in physical space (spacing $\Delta x = |\Omega|/N_x$).
The discrete $L^1$ norm is defined as $\norm{\mathbf{u}}_{L^1} = \Delta x \sum_{i=1}^{N_x} |\mathbf{u}(x_i)|$, which approximates the continuous $L^1$ norm on physical space (see the norm equivalence claim in Stage~4).
\end{remark}

\subsubsection{3.1 Kruzhkov Entropy Condition and \texorpdfstring{$L^1$}{L1} Contraction}

Consider the target hyperbolic conservation law (on physical space $\Omega \times [0,T]$):
\begin{equation}
  \partial_t \mathbf{u} + \partial_x \flux(\mathbf{u}) = 0, \qquad x \in \Omega \subset \R, \quad t \in [0,T],
\end{equation}
where $\flux$ is a flux function (e.g., Burgers flux $\flux(u) = u^2/2$), with initial condition $\mathbf{u}(\cdot,0) = \initdat$.
The Kruzhkov entropy solution $\physsol$ satisfies: for every constant $k \in \R$, in the sense of distributions,
\begin{equation}
  \partial_t |\physsol - k| + \partial_x \bigl(\mathrm{sgn}(\physsol - k)(\flux(\physsol) - \flux(k))\bigr) \leq 0.
\end{equation}

\textbf{Corollary of Kruzhkov's theorem ($L^1$-contractive semigroup).}
If $\mathbf{a}(x,t)$ and $\mathbf{b}(x,t)$ both satisfy the Kruzhkov entropy condition, then
\begin{equation}
  \norm{\mathbf{a}(\cdot,t) - \mathbf{b}(\cdot,t)}_{L^1(\Omega)} \leq \norm{\mathbf{a}(\cdot,0) - \mathbf{b}(\cdot,0)}_{L^1(\Omega)}.
\end{equation}
Set $\physsol(x) \coloneqq \physsol(x, T_{\mathrm{phys}})$ (at fixed physical time $T_{\mathrm{phys}}$; omitted in the sequel).

\subsubsection{3.2 Loss Design and BV Compactness}

The loss function $\mathcal{L}_3$ augments the standard denoising score matching loss with a BV regularizer:
\begin{equation}
  \mathcal{L}_3 = \Ldsm + \lambda_{\mathrm{BV}} \, \TV(\Dth).
\end{equation}

By Assumption~A0, together with the hard-coded $\tanh$ structure of parameterization~\eqref{eq:bv-aware}, this regularizer guarantees that the generated samples almost surely belong to the bounded BV ball $\mathcal{K}_M$.
By Helly's selection theorem, $\mathcal{K}_M$ is relatively compact in the $L^1(\Omega)$ topology.

\subsubsection{3.3 Viscous Limit (Application of Kuznetsov's Lemma)}

The reverse stage of the diffusion model, under viscosity matching $\sigtau^2(\tau) = 2\physvis \tau$, effectively introduces artificial viscosity $\physvis$. Denote by $\physsolvis$ the solution of the viscous physical PDE at $t = T_{\mathrm{phys}}$:
\begin{equation}
  \partial_t \physsolvis + \partial_x \flux(\physsolvis) = \physvis \, \partial_{xx} \physsolvis.
\end{equation}

By Assumption~A6 (based on Kuznetsov 1976):
\begin{equation}
  \norm{\physsolvis - \physsol}_{L^1(\Omega)} \leq C_K \sqrt{\physvis}.
\end{equation}

Together with A5 ($\physvis = \Theta(\varepsilon)$), we obtain
\begin{equation}\label{eq:3.1}
  \norm{\physsolvis - \physsol}_{L^1(\Omega)} \leq C_K \sqrt{C\varepsilon} = C_K \sqrt{C} \cdot \varepsilon^{1/2}. \tag{3.1}
\end{equation}

% ============================================================================
% Stage 4: Assembly of the Wasserstein Error and Emergence of O(ε^{1/2})
% ============================================================================
\subsection{Stage 4: Assembly of the Wasserstein Error and Emergence of \texorpdfstring{$\mathcal{O}(\varepsilon^{1/2})$}{O(epsilon\^{}1/2)}}\label{sec:stage4}

\subsubsection{4.1 Grid Renormalization, Norm Equivalence, and Neural Generation Error}

\textbf{Derivation of the discrete $L^1$--normalized Euclidean norm equivalence.}
The physical domain $\Omega$ is discretized into $d = N_x$ uniform grid points with spacing $\Delta x = |\Omega|/d$.
For a state vector $u \in \R^d$, the discrete $L^1$ norm is
\begin{equation}
  \norm{\mathbf{u}}_{L^1} = \Delta x \sum_{i=1}^d |u_i|.
\end{equation}

From the $\ell^1$--$\ell^2$ norm inequality ($\sum |u_i| \leq \sqrt{d} \sqrt{\sum u_i^2} = \sqrt{d}\norm{u}$), substituting $\Delta x$, we obtain
\begin{equation}\label{eq:4.1}
  \norm{\mathbf{u}}_{L^1} \leq \Delta x \sqrt{d} \; \norm{u} = \sqrt{|\Omega|} \cdot \sqrt{\Delta x} \; \norm{u}. \tag{4.1}
\end{equation}

\textbf{Deriving the $L^1$ bound for the neural generation error.}
Note that in the renormalized Assumption~A3, the error bound is given in the form $\E[\Delta x \norm{\errsc}^2] \leq \varepsilon^2$, which means that for the pure Euclidean norm, we have $\E[\norm{\errsc}^2] \leq \varepsilon^2 / \Delta x$.

Recall from Stage~2, the Gronwall integration and trajectory error bound Lemma~2.2 give, for the standard Euclidean distance $E(0)$:
\begin{equation}\label{eq:4.2}
  \E_z\bigl[\norm{\uhat(0) - \unat(0)}\bigr] \leq C_1 \; \E_z\bigl[\norm{\errsc}\bigr] \leq C_1 \frac{\varepsilon}{\sqrt{\Delta x}}. \tag{4.2}
\end{equation}

The corresponding difference of physical functions is $\uth - \physsolvis$. Combining with the norm scaling relation~(4.1):
\begin{equation}
  \E_z\bigl[\norm{\uth - \physsolvis}_{L^1(\Omega)}\bigr] \leq \sqrt{|\Omega|} \sqrt{\Delta x} \cdot \E_z\bigl[\norm{\uhat(0) - \unat(0)}\bigr] \leq \sqrt{|\Omega|} \sqrt{\Delta x} \cdot C_1 \frac{\varepsilon}{\sqrt{\Delta x}}.
\end{equation}

Here $d$ and $\Delta x$ are \emph{perfectly canceled}! Regardless of how fine the grid is, the error bound remains valid, yielding a grid-independent continuum limit bound:
\begin{equation}\label{eq:4.3}
  \E_z\bigl[\norm{\uth - \physsolvis}_{L^1(\Omega)}\bigr] \leq C_1 \sqrt{|\Omega|} \, \varepsilon \eqqcolon C_2 \varepsilon. \tag{4.3}
\end{equation}

\begin{remark}[Place of the Gagliardo--Nirenberg interpolation]
An earlier draft used a Gagliardo--Nirenberg-type interpolation inequality $\norm{h}_{L^1} \leq C_{\mathrm{GN}} \norm{h}_{L^2}^{2/3} \TV(h)^{1/3}$ to convert $L^2$ error into $L^1$ error, obtaining an $\mathcal{O}(\varepsilon^{2/3})$ bound.
The present revised proof uses normalized $L^2$-matching combined with the norm equivalence~(4.1), obtaining a sharper, dimension-independent $L^1$ error $\mathcal{O}(\varepsilon)$ directly from the scaled Euclidean error.
\end{remark}

\subsubsection{4.2 Triangle Inequality Decomposition of the Wasserstein Distance}

The 1-Wasserstein distance (Kantorovich--Rubinstein dual):
\begin{equation}
  \Wass{1}(\mu, \nu) = \inf_{\pi \in \Pi(\mu,\nu)} \int \norm{\mathbf{u} - \mathbf{v}}_{L^1(\Omega)} \, d\pi(\mathbf{u},\mathbf{v}).
\end{equation}

Let $\delta_{\physsolvis}$ be the Dirac measure at the viscous solution (given the initial data $\initdat$, the viscous PDE solution is unique, hence deterministic). Decompose via the triangle inequality:
\begin{equation}\label{eq:4.4}
  \Wass{1}(\muth, \delta_{\physsol})
  \leq \underbrace{\Wass{1}(\muth, \delta_{\physsolvis})}_{\text{Neural generation error (Term~A)}}
  + \underbrace{\Wass{1}(\delta_{\physsolvis}, \delta_{\physsol})}_{\text{Viscous limit error (Term~B)}}. \tag{4.4}
\end{equation}

\textbf{Term~A (Neural generation error).}
Since $\uhat(0)$ and $\unat(0)$ are deterministically generated from the same prior noise $z$ (the ODE has a unique solution for fixed $z$), there exists a natural pointwise coupling:
\begin{equation}
  \pi_z = \bigl(\uth(z),\; \physsolvis\bigr)_\# \, p_{T_d},
\end{equation}
where $p_{T_d}$ is the prior distribution and $\physsolvis$ is regarded as a constant vector (independent of $z$). By the feasibility of this coupling and~(4.3):
\begin{equation}\label{eq:4.5}
  \Wass{1}(\muth, \delta_{\physsolvis})
  \leq \int \norm{\uth(z) - \physsolvis}_{L^1(\Omega)} \, d p_{T_d}(z)
  = \E_z\bigl[\norm{\uth - \physsolvis}_{L^1(\Omega)}\bigr]
  \leq C_2 \varepsilon. \tag{4.5}
\end{equation}

\textbf{Term~B (Viscous limit error).}
Both $\delta_{\physsolvis}$ and $\delta_{\physsol}$ are single-point distributions; $\Wass{1}$ degenerates to the $L^1$ distance:
\begin{equation}
  \Wass{1}(\delta_{\physsolvis}, \delta_{\physsol}) = \norm{\physsolvis - \physsol}_{L^1(\Omega)}.
\end{equation}

From~(3.1) (the joint conclusion of A6 and A5):
\begin{equation}\label{eq:4.6}
  \Wass{1}(\delta_{\physsolvis}, \delta_{\physsol}) \leq C_K \sqrt{C} \cdot \varepsilon^{1/2}. \tag{4.6}
\end{equation}

\subsubsection{4.3 Final Conclusion}

Substituting~(4.5) and~(4.6) into~(4.4):
\begin{equation}
  \Wass{1}(\muth, \delta_{\physsol}) \leq C_2 \varepsilon + C_K \sqrt{C} \cdot \varepsilon^{1/2}.
\end{equation}

For $\varepsilon \in (0,1)$, we have $\varepsilon \leq \varepsilon^{1/2}$, so the higher-order term $\mathcal{O}(\varepsilon)$ is absorbed by the lower-order term $\mathcal{O}(\varepsilon^{1/2})$:
\begin{equation}
  \Wass{1}(\muth, \delta_{\physsol}) \leq \bigl(C_2 + C_K \sqrt{C}\bigr) \cdot \varepsilon^{1/2} \eqqcolon C_{\mathrm{final}} \cdot \varepsilon^{1/2}.
\end{equation}

Here $C_{\mathrm{final}} = C_2 + C_K \sqrt{C}$ depends only on $M, T_d, T_{\mathrm{phys}}, \kappa_0, L_{\mathrm{sm}}, \physvis/\varepsilon$, all independent of $\varepsilon$. Thus
\begin{equation}
  \boxed{\displaystyle \Wass{1}(\muth, \delta_{\physsol}) \leq C_{\mathrm{final}} \, \varepsilon^{1/2}.}
\end{equation}

This completes the proof of Theorem~\ref{thm:improved-rate}.
\end{proof}

% ============================================================================
% Appendix A: Symbol Correspondence Table (Old → Revised)
% ============================================================================
\subsection*{Appendix A: Symbol Correspondence Table (Old \texorpdfstring{$\to$}{to} Revised)}

\begin{table}[h]
\centering
\caption{Symbol correspondence table}
\begin{tabular}{lll}
\toprule
Old symbol & Revised symbol & Reason \\
\midrule
$u \in \R^d$ (state point) and $u(x,t)$ (PDE solution) confused & State space: $u \in \R^d$; physical space: $\mathbf{u}(x)$ & Disambiguate function/variable \\
$\uth$ used for both true trajectory and PDE solution & $\unat(\tau)$ (true trajectory); $\physsol$ (PDE solution) & Distinguish two types of ``truth'' \\
$\uhat$ confused as trajectory/sample/distribution & $\uhat(\tau)$ (trajectory); $\uth$ (sample); $\muth$ (distribution) & Clarify semantics at different stages \\
$\nabla_u$ vs $\dx$ confused & $\nabla_u$ (state-space gradient); $\dx$ (physical space partial derivative) & Distinguish two different spaces \\
$\sup_u \lambda_{\max}$ with ambiguous scope & $\sup_{u \in \R^d} \lambda_{\max}$ & Clarify $u$ as a state-space point \\
$\norm{E(\tau)}$ with ambiguous norm type & $E(\tau) \coloneqq \norm{\uhat(\tau) - \unat(\tau)}$ ($\R^d$ Euclidean norm) & Clarify as state-space norm \\
Dual role of $u$ in residual field $\errsc(u,\tau)$ & The parameter $u$ in $\errsc(u,\tau)$ is a state-space point & Retain $u$ as parameter name, add scope annotation \\
Relation between times $\tau$ and $t$ unclear & Time mapping paragraph explicitly defined & Explicitly declare orthogonality and physical meaning \\
\bottomrule
\end{tabular}
\end{table}

% ============================================================================
% Appendix B: Summary of Modifications in this Revision
% ============================================================================
\subsection*{Appendix B: Summary of Modifications}

\begin{table}[h]
\centering
\caption{Summary of modifications}
\begin{tabular}{llc}
\toprule
No. & Modification & Treatment \\
\midrule
M1 & Rigorous gradient bound claim in Corollary~1.4 & Downgraded to heuristic remark in \S\ref{sec:stage1}; not relied upon subsequently \\
M2 & The $\int \norm{\errsc} \, d\tau$ bound in Lemma~2.2 & Directly from A3 + Cauchy--Schwarz; no dependence on Corollary~1.4 \\
M3 & GN interpolation in Stage~4 & Removed; replaced with discrete norm equivalence (4.2) \\
M4 & Viscous error Term~B & Directly cites A5 + A6 (Equation~(3.1)); logical chain closed \\
M5 & Theorem statement & Explicitly references A0--A6; no dangling assumptions \\
M6 & Assumption system & Newly added A0--A6, with verifiability annotations and engineering approximation labels \\
\bottomrule
\end{tabular}
\end{table}

% ============================================================================
% Appendix C: Technical Points Requiring Further Rigorization
% ============================================================================
\subsection*{Appendix C: Technical Points Requiring Further Rigorization and Scope of Engineering Approximations}

\begin{table}[h]
\centering
\caption{Checklist of technical points}
\small
\begin{tabular}{llcl}
\toprule
No. & Technical point & Status & Scope of influence \\
\midrule
C1.1-P & ``Finitely separated Gaussian clusters'' assumption of Lemma~1.1 & Upgraded to A1 & If A1 fails, the score singularity decomposition may be inaccurate \\
C1.4-G & Residual gradient bound of Corollary~1.4 & Left to future work & This proof does not rely on this bound \\
C2.1-$\kappa$ & Training-time enforcement of $\jumpamp > 0$ & Upgraded to A2 & If not strictly enforced, conclusion degrades to $\mu(\tau) = \mathcal{O}(1)$ \\
EA1 & Applicability of A1 to general continuous distributions & Engineering approximation & Affects the uniformity of the $\mathcal{O}(\tau)$ remainder in Lemma~1.1 \\
EA2 & Feasibility of schedule $\physvis = \Theta(\varepsilon)$ in A5 & Engineering approximation & If $\physvis \gg \varepsilon$, Term~B dominates; if $\physvis \ll \varepsilon$, Term~A dominates \\
P4.1-GN & GN interpolation argument in the continuum limit & Retained in remark & Does not affect conclusion on fixed grid \\
K3.3-$\nu$ & Generalization of Kuznetsov bound in A6 & Rigorously established & Only requires initial BV bound (guaranteed by A0) \\
\bottomrule
\end{tabular}
\end{table}

% ----- A1.4 · Theorem 4 · Shock-location admissibility -----
\subsection{Proof of Theorem~\ref{thm:shock-location} (Shock-location admissibility)}
\label{appx:proofs:thm4}
%==============================================================================
% Appendix: Proof of Theorem 4 — Shock-Location Admissibility
% Auto-converted from Docs/proof/Theorem 4 draft.md
%==============================================================================
\begin{theorem}[Shock-location admissibility]
\label{thm:shock-location}
In the limit $\sigtau \to 0$ (i.e., the final step of inference), the solution produced by the
proposed sampler satisfies the Rankine--Hugoniot condition at each shock location~$x_{s}(t)$
almost everywhere,
\begin{equation}
\label{eq:rh-appx}
\dot x_{s}(t) = \frac{\flux(u_{L}) - \flux(u_{R})}{u_{L} - u_{R}},
\end{equation}
and the Lax entropy condition $\flux'(u_{L}) \ge \dot x_{s} \ge \flux'(u_{R})$.
\end{theorem}

\begin{proof}

\noindent\textbf{Assumptions.}

\begin{enumerate}
\item \textbf{(Perfect score matching at the interfacial layer.)}
  The BV-aware score parameterization~\eqref{eq:bv-aware} is optimized to global optimality
  in a neighborhood of the discontinuity, i.e.\ $\sth(u,\tau) \equiv \nabla_u \log \ptau(u)$.
  We thus analyze the true score field.

\item \textbf{(Single-shock target manifold.)}
  The target distribution $\rhotphys(u)$ concentrates on a submanifold $\mathcal{M}_{t}$
  in function space (or its discretization~$\R^{N}$).%
  For the sampling trajectory under consideration, the corresponding target
  Kruzhkov entropy solution $\physsol(x,t)$ at the current time~$t$ possesses a single
  isolated shock located at~$x_{s}(t)$, with left and right limits
  \begin{equation}
  u_{L} \coloneqq \physsol(x_{s}^{-},t), \qquad
  u_{R} \coloneqq \physsol(x_{s}^{+},t).
  \end{equation}

\item \textbf{(Local planarity in high dimension.)}
  Since the grid dimension~$N$ is large, in the limit $\tau\to0$ the separating
  manifold of data modes (i.e., the locus where the score exhibits a shock) can be
  locally approximated as a hyperplane.
  Its signed distance function~$\phissh(u,\tau)$ satisfies
  $\abs{\nabla_u \phissh} = 1$.

\item \textbf{(Compressive score gradient.)}
  When the reverse sampling trajectory $u(\tau)$ approaches the target
  manifold~$\mathcal{M}_{t}$, the local flow induced by the network
  parameterization~$\sth$ is compressive along the shock
  normal~$\nabla\phissh$:
  \begin{equation}
  \partial_{\phissh} \sth \cdot \nabla\phissh < 0.
  \end{equation}
  This guarantees that the score field acts as a centripetal force, pushing
  probability mass toward the interfacial manifold.
\end{enumerate}

%==========================================================================
% Part I: Geometric Correspondence via the Viscosity-Matched Schedule
%==========================================================================
\medskip
\noindent\textbf{Part I: Geometric correspondence via the viscosity-matched schedule.}

Let the forward noising process be a variance-exploding SDE.
Under the schedule~\eqref{eq:visc-matched}, the noise variance is tied to the physical
viscosity~$\physvis$:
\begin{equation}
\sigtau^{2} = 2\physvis\,\tau, \qquad \tau \in [0,T_{d}].
\end{equation}
The forward SDE then reads
\begin{equation}
du = \sqrt{2\physvis}\; dw.
\end{equation}
The corresponding Fokker--Planck equation for the density $p(u,\tau)$ in data space~$\R^{N}$ is
\begin{equation}
\label{eq:fp}
\dtau p(u,\tau) = \physvis\,\Lap_{u} p(u,\tau).
\end{equation}

Differentiating the score $\score(u,\tau) = \nabla_u \log p = \nabla_u p / p$ with
respect to~$\tau$ and invoking~\eqref{eq:fp} gives the exact evolution of the score:
\begin{align}
\dtau \score
  &= \physvis\,\nabla_u\!\Bigl( \frac{\Lap_u p}{p} \Bigr) \notag\\
  &= \physvis\,\Lap_u \score + \physvis\,\nabla_u(\abs{\score}^{2}).
\end{align}
Since $\score$ is a gradient field, $\nabla_u(\abs{\score}^{2}) = 2(\score\cdot\nabla_u)\score$.
Note that the relation $\score = \nabla_u\log p$ is essentially the classical
\emph{inverse Cole--Hopf transform}, which explicitly converts the linear heat equation
(Fokker--Planck) into a nonlinear PDE.
We thus obtain the \textbf{high-dimensional viscous Burgers equation at the score level}:
\begin{equation}
\label{eq:score-burg}
\dtau \score + (-2\physvis\,\score\cdot\nabla_u)\,\score = \physvis\,\Lap_u \score.
\end{equation}

\begin{lemma}[Effective viscosity equivalence]
\label{lem:visc-equiv}
Via the schedule~\eqref{eq:visc-matched}, the effective viscosity of the score dynamics is
numerically identical to the physical viscosity~$\physvis$ of the target equation.
\end{lemma}

The target PDE with small physical viscosity is
\begin{equation}
\label{eq:target-pde}
\dt u(x,t) + \dx \flux(u(x,t)) = \physvis\,\partial_{xx} u(x,t).
\end{equation}
Under discretization, the components of $u \in \R^{N}$ satisfy $u_i \approx u(x_i,t)$.
The target distribution $\rhotphys(u)$ (the boundary condition at $\tau=0$) is the physical
weak-solution distribution.
In the limit $\tau\to0$, the reverse sampling trajectory $u(\tau)$ is pulled by the
score~$\score(u,\tau)$ toward the support manifold~$\mathcal{M}_{t}$ of~$\rhotphys(u)$.

The BV-aware score parameterization~\eqref{eq:bv-aware} reads
\begin{equation}
\sth(u,\tau)
  = \nabla\phisbg + \frac{\jumpamp}{2}\tanh\!\Bigl(\frac{\phissh}{2}\Bigr)\nabla\phissh.
\end{equation}
Near the interface, the smooth background $\nabla\phisbg$ is negligible, and the true
score exhibits the asymptotic inner-scale form
\begin{equation}
\label{eq:score-inner}
\score(u,\tau)
  \sim \frac{\jumpamp(u)}{2}
       \tanh\!\Bigl( \frac{\operatorname{dist}(u,\mathcal{M}_{t})}
                          {\sqrt{2\physvis\,\tau}} \Bigr)
       \mathbf{n}(u),
\end{equation}
where $\mathbf{n}(u) = \nabla_u \operatorname{dist}(u,\mathcal{M}_{t})$ is the normal vector
to the manifold.

\begin{lemma}[Geometric correspondence in the double limit]
\label{lem:geom-corresp}
When $\tau\to0$ and $\physvis\to0$ hold simultaneously, the $\tanh$ profile
in~\eqref{eq:score-inner} degenerates to a step function.
Consequently, the location where the score field exhibits a shock is determined in data space
by $\operatorname{dist}(u,\mathcal{M}_{t}) = 0$.
Mapped back to physical space, this manifold boundary coincides exactly with the grid
coordinates where $u(x,t)$ exhibits a jump, i.e., the physical shock position~$x_{s}(t)$.
\end{lemma}

%==========================================================================
% Part II: Weak-Form Limit and Rankine--Hugoniot Condition
%==========================================================================
\medskip
\noindent\textbf{Part II: Weak-form limit and derivation of the Rankine--Hugoniot condition.}

\smallskip
\noindent\textit{2.1  Tweedie profile near the shock.}

For the reverse diffusion, the clean physical state $\uhat_{0}(x,t)$ is estimated via the
Tweedie formula
\begin{equation}
\uhat_{0}(x,t) = u(\tau) + \sigtau^{2}\,\score(u,\tau).
\end{equation}
Under the schedule~\eqref{eq:visc-matched} and the BV-aware parameterization, restricted
to a one-dimensional physical section, the signed distance expands as
$\displaystyle \phissh \approx \frac{x - x_{s}(t)}{\sqrt{2\physvis\,\tau}}$.

\begin{lemma}[Viscous shock profile]
\label{lem:visc-profile}
In a neighborhood of $x_{s}(t)$, the Tweedie-reconstructed physical field
$\uhat^{\tau}(x,t)$ possesses a smooth, monotone shock transition layer.
To make the physical scales explicit, define the layer thickness
$\varepsilon = \sigtau = \sqrt{2\physvis\,\tau}$.
It is described by
\begin{equation}
\label{eq:visc-profile}
\uhat^{\tau}(x,t)
  \approx \frac{u_{L}+u_{R}}{2}
        - \frac{u_{L}-u_{R}}{2}
          \tanh\!\Bigl( \frac{x - x_{s}(t)}{2\varepsilon} \Bigr).
\end{equation}
Thus $\uhat^{\tau}$ is strictly between $u_{L}$ and $u_{R}$ and has controlled total variation.
\end{lemma}

\smallskip
\noindent\textit{2.2  Godunov guidance and weak form.}

During generation, an entropy-compatible Godunov flux difference is used as PDE guidance.
Thus, as $\tau\to0$, the generated sequence $\uhat^{\tau}$ approximately satisfies a weak form:
for every smooth test function $\varphi \in C_{c}^{\infty}(\R\times(0,T))$,
\begin{equation}
\label{eq:weak-form}
I(\varepsilon)
  = \iint_{\R\times\R^{+}}
    \bigl( \uhat^{\tau}\,\partial_{t}\varphi + \flux(\uhat^{\tau})\,\partial_{x}\varphi \bigr)
    \,dx\,dt
  = \mathcal{O}(\varepsilon).
\end{equation}

Because the $\tanh$ profile guarantees $\uhat^{\tau} \in L^{\infty} \cap \BV$, the dominated
convergence theorem allows the limit $\varepsilon\to0$ to be taken inside the integral.
\textbf{This step is particularly emphasized: the local monotonicity induced by
parameterization~\eqref{eq:bv-aware} is the key topological constraint that renders this weak
limit legitimate.}
Traditional diffusion models (e.g., those forcing a Gaussian fit across shock interfaces)
inevitably produce Gibbs oscillations, causing the~$\BV$ norm to blow up in the limit and
precluding application of dominated convergence.
The profile~\eqref{eq:visc-profile} converges strongly to the piecewise-constant function
\begin{equation}
\physsol(x,t) =
\begin{cases}
u_{L}, & x < x_{s}(t),\\[2pt]
u_{R}, & x > x_{s}(t).
\end{cases}
\end{equation}

Partition the spacetime domain along the shock trajectory $\Gamma\!: x = x_{s}(t)$ into a
left region $D^{-}$ ($x < x_{s}$) and a right region $D^{+}$ ($x > x_{s}$).
Then
\begin{equation}
\label{eq:limit-split}
0 = \lim_{\varepsilon\to0} I(\varepsilon)
  = \iint_{D^{-}} \bigl( u_{L}\,\partial_{t}\varphi + \flux(u_{L})\,\partial_{x}\varphi \bigr)
    \,dx\,dt
  + \iint_{D^{+}} \bigl( u_{R}\,\partial_{t}\varphi + \flux(u_{R})\,\partial_{x}\varphi \bigr)
    \,dx\,dt.
\end{equation}
Inside each $D^{\pm}$, the solution $u_{L},u_{R}$ is a strong solution of the conservation law.
Applying Green's theorem to transfer derivatives from $\varphi$ onto~$u$,
\begin{equation}
\label{eq:green}
\begin{aligned}
\iint_D \bigl( u\,\partial_{t}\varphi + \flux(u)\,\partial_{x}\varphi \bigr) \,dx\,dt
&= \int_{\partial D}
   \varphi\,\bigl( u\,n_{t} + \flux(u)\,n_{x} \bigr)\,ds \\
&\quad - \iint_D \varphi\,\bigl( \partial_{t}u + \partial_{x}\flux(u) \bigr) \,dx\,dt.
\end{aligned}
\end{equation}
Since the PDE holds pointwise in $D^{\pm}$, the area term vanishes.
Because $\varphi$ has compact support, only the boundaries along~$\Gamma$ contribute.
On $\Gamma$, the outward normal (pointing toward $D^{+}$) is
\begin{equation}
\mathbf{n} = (n_{x}, n_{t})
           = \frac{1}{\sqrt{1 + \dot x_{s}^{2}}}
             \bigl( 1,\, -\dot x_{s} \bigr).
\end{equation}
Accounting for orientation, the jump condition emerges:
\begin{equation}
\label{eq:jump-cond}
\int_{\Gamma} \varphi \cdot
\Bigl[
  \bigl( u_{L}(-\dot x_{s}) + \flux(u_{L}) \bigr)
  - \bigl( u_{R}(-\dot x_{s}) + \flux(u_{R}) \bigr)
\Bigr]
\frac{1}{\sqrt{1 + \dot x_{s}^{2}}}
\,ds = 0.
\end{equation}
Since $\varphi$ is arbitrary along $\Gamma$, the algebraic bracket must vanish identically,
yielding
\begin{equation}
\dot x_{s}\,(u_{L} - u_{R}) = \flux(u_{L}) - \flux(u_{R}).
\end{equation}
For a genuine shock $u_{L} \neq u_{R}$, division by the difference gives the
Rankine--Hugoniot condition
\begin{equation}
\label{eq:rh-derived}
\dot x_{s}(t) = \frac{\flux(u_{L}) - \flux(u_{R})}{u_{L} - u_{R}}.
\qquad\square
\end{equation}

%==========================================================================
% Part III: Entropy Compatibility and the Lax Condition
%==========================================================================
\medskip
\noindent\textbf{Part III: Entropy compatibility and the Lax condition.}

\smallskip
\noindent\textit{Background: the non-uniqueness of weak solutions.}
The Rankine--Hugoniot condition derived above is a \emph{necessary} but not sufficient
condition for a weak solution.
Physically inadmissible ``expansion shocks'', which dissipate energy in an
entropy-violating manner, also satisfy the R--H condition in nonlinear evolution models.
Only by imposing a rigorous entropy condition---the Lax entropy condition, motivated by
physical mechanisms such as the Double-Burgers coupling---do we ensure that the generated
solution is not merely a ``numerical fit'' to data, but a genuine approximation of the
\emph{physically correct} conservation law.

\smallskip
\noindent\textit{3.1  Divergence--characteristic relation.}

For a convex flux $\flux$ (e.g., the Burgers flux $\flux(u)=u^{2}/2$), a physical shock
must satisfy the Lax entropy condition
\begin{equation}
\label{eq:lax}
\flux'(u_{L}) > \dot x_{s} > \flux'(u_{R}).
\end{equation}
In the reverse-time generative ODE,
\begin{equation}
\frac{du}{d\tau} = -\frac12\,\beta(\tau)\bigl[ u + \sth(u,\tau) \bigr],
\end{equation}
Tweedie reconstruction yields
$\partial_{x}\uhat \approx \partial_{x}(u + \sigtau^{2}\,\sth) \to -\infty$
near $x_{s}$.
This singular negative gradient corresponds to the compression of characteristics
toward the shock in physical space.

\smallskip
\noindent\textit{3.2  Kruzhkov entropy inequality.}

To prove that $\uhat$ is an entropy solution, we verify the Kruzhkov inequality for every
constant~$k$ (in the sense of distributions):
\begin{equation}
\label{eq:kruz}
\partial_{t}\,\abs{\uhat - k}
  + \partial_{x}\bigl[ \operatorname{sgn}(\uhat - k)\,(\flux(\uhat) - \flux(k)) \bigr]
  \le 0.
\end{equation}
By Theorem~1 (Double-Burgers coupling), the score field itself satisfies a viscous
Burgers equation.
Hence the generative dynamics driven by $\sth$ constitute a JKO gradient flow that
minimizes an energy functional.
A jump $u_{L} \to u_{R}$ violating the Lax condition (i.e., $u_{L} < u_{R}$) would
correspond to an expansion (rarefaction) fan whose characteristics diverge outward.
Under the BV-aware parameterization~\eqref{eq:bv-aware}, the sign of the score gradient is
determined by the $\tanh$ profile.

\smallskip
\noindent\textit{3.3  Proof by contradiction.}

Suppose a non-physical shock (an expansion shock) satisfies the R--H condition but has
$u_{L} < u_{R}$.

\begin{enumerate}
\item[\textbf{(i)}] \textbf{Score behavior.}
  The Tweedie formula then requires the score~$\sth$ to supply an outward ``push'' in order
  to sustain this unstable discontinuity.

\item[\textbf{(ii)}] \textbf{Burgers consistency loss.}
  The loss term $\Lburg$ (Burgers consistency) enforces
  $\dtau \score + (\score\cdot\nabla)\score = \physvis\,\Lap\score$.
  An expansion jump would be immediately regularized by the diffusion term~$\Lap\score$,
  producing a large residual in~$\Lburg$.

\item[\textbf{(iii)}] \textbf{Conclusion.}
  Minimizing the total loss systematically eliminates all non-physical jumps, because only
  compressive configurations satisfying the Lax condition~\eqref{eq:lax} are compatible with
  the stable evolution operator of~$\Lburg$.
\end{enumerate}

\smallskip
\noindent\textit{3.4  Stability in the $\tau\to0$ limit.}

Finally, the Oleinik entropy condition provides the estimate near the shock:
\begin{equation}
\frac{\uhat(x+a,t) - \uhat(x,t)}{a} \le \frac{C}{t},
\end{equation}
which is guaranteed by the controlled width of the $\tanh$ layer in
parameterization~\eqref{eq:bv-aware}.
As $\tau$ decreases, the profile steepens but remains strictly monotone whenever $u_{L}>u_{R}$,
thereby precluding non-physical oscillations.

\medskip
\noindent
Combining the three parts above, the generated solution satisfies both the Rankine--Hugoniot
condition and the Lax entropy condition almost everywhere.
\end{proof}


% ----- A1.5 · Theorem 5 · JKO correspondence -----
\subsection{Proof of Theorem~\ref{thm:jko} (JKO correspondence)}
\label{appx:proofs:thm5}
%==============================================================================
% Appendix: Proof of Theorem 5 — EntroDiff Reverse Sampler as JKO Gradient Flow
% Auto-converted from Docs/proof/Theorem 5 draft.md
%==============================================================================
\begin{theorem}[JKO correspondence]
\label{appx-thm:jko}
Let $\ptau$ denote the density of the reverse sampler at diffusion time
$\tau \in [0,T]$, and let $\rhotrue$ be the target distribution induced by the PDE
solution in data space.
Define the free energy functional
\begin{equation}
\mathcal{F}(\rho) = \Ent(\rho \mid \rhotrue) + \lambda\,\mathcal{R}(\rho),
\qquad
\Ent(\rho \mid \rhotrue) = \int_{\R^{N}} \rho(x)\,\log\frac{\rho(x)}{\rhotrue(x)}\,dx,
\end{equation}
where $\mathcal{R}$ is a density functional encoding the PDE constraint and $\lambda \ge 0$.

Assume the following hypotheses:
\begin{enumerate}
\item \textbf{(Exact score limit.)}
  After sufficient training, $\sth(x,\tau) = \nabla_x \log \ptau(x)$ pointwise.

\item \textbf{(Viscosity-matched schedule.)}
  $a(\tau) = g(\tau)^{2}/2 > 0$, aligned with the physical viscosity.

\item \textbf{(PDE constraint as a differentiable functional.)}
  $\mathcal{R}$ is Fr\'{e}chet-differentiable on $\mathcal{P}_{2}(\R^{N})$, with
  first variation $\delta \mathcal{R} / \delta \rho$ well-defined;
  $\nabla_x (\delta \mathcal{R} / \delta \rho)$ is locally bounded.

\item \textbf{(Rigid $\Lburg$ consistency.)}
  The Burgers consistency penalty~$\Lburg$ is sufficiently strong that, in the
  infinite-capacity noiseless limit, the learned score field satisfies the viscous
  Burgers structure
  \begin{equation}
  \dtau \score_{\tau} + 2\,\score_{\tau}\cdot\nabla\,\score_{\tau}
    = \Lap \score_{\tau}.
  \end{equation}
  Moreover, $\ptau \in \mathcal{P}_{2}(\R^{N})$ is smooth in $(x,\tau)$ and strictly
  positive ($\ptau > 0$).

\item \textbf{(PDE guidance as a gradient drift.)}
  In the continuous-time limit, the combined PDE constraint losses $\Lent$
  (Kruzhkov entropy penalty) and $\Lburg$ (Burgers consistency) induce an additional
  guidance velocity $g_{\tau}(x)$ in the reverse ODE that converges to the spatial
  gradient of a functional.
  Specifically, the total effective velocity field of the EntroDiff reverse sampler is
  \begin{equation}
  \label{eq:total-vel-stmt}
  v_{\tau}(x)
    = -a(\tau)\Bigl[
        \nabla_x\log\ptau(x)
        + \nabla_x V(x)
        + \lambda\,\nabla_x\frac{\delta\mathcal{R}}{\delta\rho}(\ptau)(x)
      \Bigr],
  \end{equation}
  where $V$ is the potential of the target distribution ($\rhotrue \propto e^{-V}$).
  The guidance contributions $\nabla V$ and $\lambda\nabla(\delta\mathcal{R}/\delta\rho)$
  arise from the gradient of the PDE losses with respect to the sample trajectory,
  which by design converge to these functional-gradient limits when the losses are driven
  to zero in training.
\end{enumerate}

Then the density evolution of the reverse sampler satisfies
\begin{equation}
\label{eq:density-flow}
\dtau \ptau = \nabla \cdot \Bigl( a(\tau)\, \ptau\, \nabla\frac{\delta\mathcal{F}}{\delta\rho}(\ptau) \Bigr).
\end{equation}
Equivalently, under the time reparameterization
$\theta(\tau) = \int_{0}^{\tau} a(r)\,dr$, we have
\begin{equation}
\label{eq:W2-flow}
\partial_{\theta} p_{\theta}
  = \nabla \cdot \Bigl( p_{\theta}\, \nabla\frac{\delta\mathcal{F}}{\delta p}(p_{\theta}) \Bigr),
\end{equation}
which is precisely the $\Wass{2}$-gradient flow of~$\mathcal{F}$.
Consequently, by the JKO theory, a single time-discrete step takes the variational form
\begin{equation}
\label{eq:jko-step}
p^{n+1}
  = \arg\min_{p \in \mathcal{P}_{2}(\R^{N})}
    \Bigl\{ \mathcal{F}(p) + \frac{1}{2\Delta\theta}\,\Wass{2}^{2}(p, p^{n}) \Bigr\}.
\end{equation}
\end{theorem}

\begin{proof}

%==========================================================================
% Step 1: Reverse Sampler as a Continuity Equation
%==========================================================================
\medskip
\noindent\textbf{Step 1: Reverse sampler as a continuity equation.}

Let $X_{\tau}$ denote the particle trajectory of the reverse sampler.
If its continuous-time limit satisfies the ODE
\begin{equation}
\frac{d X_{\tau}}{d\tau} = v_{\tau}(X_{\tau}),
\end{equation}
then by the standard pushforward formula, the particle distribution
$\ptau = \Law(X_{\tau})$ satisfies the continuity equation
\begin{equation}
\label{eq:continuity}
\dtau \ptau + \nabla \cdot (\ptau\,v_{\tau}) = 0.
\end{equation}
Thus, identifying the velocity field $v_{\tau}$ is the central task.

For any test function $\varphi \in C_{c}^{\infty}(\R^{N})$, the identity follows from the
chain rule:
\begin{align}
\frac{d}{d\tau} \int_{\R^{N}} \varphi(x)\,\ptau(x)\,dx
  &= \frac{d}{d\tau}\,\E[\varphi(X_{\tau})] \notag\\
  &= \E[\nabla\varphi(X_{\tau}) \cdot v_{\tau}(X_{\tau})] \notag\\
  &= \int \nabla\varphi(x) \cdot v_{\tau}(x)\,\ptau(x)\,dx.
\end{align}
Integration by parts of the right-hand side yields
\begin{equation}
\frac{d}{d\tau} \int \varphi\,\ptau\,dx
  = - \int \varphi(x)\,\nabla\cdot(\ptau\,v_{\tau})(x)\,dx,
\end{equation}
and since $\varphi$ is arbitrary, the continuity equation~\eqref{eq:continuity} holds in
the sense of distributions (and pointwise under the smoothness of $\ptau$ and $v_{\tau}$).

%==========================================================================
% Step 2: Score Drift and PDE Guidance — Total Effective Velocity
%==========================================================================
\medskip
\noindent\textbf{Step 2: Score drift and PDE guidance --- the total effective velocity.}

We identify $v_{\tau}$ as the sum of two contributions.

\smallskip
\noindent\textit{(a)  Score drift (from denoising score matching).}
Under Hypothesis~1 (exact score) and Hypothesis~4 (rigid $\Lburg$), the network score
converges to the true score:
\begin{equation}
\sth(x,\tau) = \nabla_x \log \ptau(x).
\end{equation}
The probability flow ODE induced by score matching alone has the velocity
\begin{equation}
\label{eq:score-vel}
v_{\tau}^{\mathrm{score}}(x) = -a(\tau)\,\sth(x,\tau) = -a(\tau)\,\nabla_x\log\ptau(x).
\end{equation}

\smallskip
\noindent\textit{(b)  PDE guidance drift (from constraint losses).}
The EntroDiff loss family includes PDE-constraint terms beyond denoising score matching:
$\Lent$ (Kruzhkov entropy penalty) and $\Lburg$ (Burgers consistency).
These losses are functionals of the generated trajectory that penalize deviations from
physical admissibility.
In the continuous-time and infinite-data limit, their gradients with respect to the sample
position induce an additional drift $g_{\tau}(x)$.

By Hypothesis~5, as the constraint losses are driven to zero, this guidance drift converges
to the gradient of the potential that defines the target distribution plus the variational
gradient of the PDE-constraint functional:
\begin{equation}
\label{eq:guidance-vel}
g_{\tau}(x) = - \Bigl( \nabla_x V(x) + \lambda\,\nabla_x\frac{\delta\mathcal{R}}{\delta\rho}(\ptau)(x) \Bigr).
\end{equation}
Here $\nabla V = -\nabla\log\rhotrue$ is the drift toward the target $\rhotrue \propto e^{-V}$,
and $\lambda\nabla(\delta\mathcal{R}/\delta\rho)$ is the PDE-constraint drift
(see Step~4 for the construction of~$\mathcal{R}$).

\smallskip
\noindent\textit{(c)  Total effective velocity.}
Summing~\eqref{eq:score-vel} and~\eqref{eq:guidance-vel}, and absorbing the common
factor $-a(\tau)$,
\begin{equation}
\label{eq:total-vel}
v_{\tau}(x)
  = -a(\tau)\Bigl[
      \nabla_x\log\ptau(x)
      + \nabla_x V(x)
      + \lambda\,\nabla_x\frac{\delta\mathcal{R}}{\delta\rho}(\ptau)(x)
    \Bigr].
\end{equation}
Substituting into the continuity equation gives
\begin{equation}
\label{eq:cont-with-vel}
\dtau \ptau
  = \nabla \cdot \Bigl(
      a(\tau)\,\ptau\,
      \bigl[ \nabla\log\ptau + \nabla V + \lambda\nabla\tfrac{\delta\mathcal{R}}{\delta\rho} \bigr]
    \Bigr).
\end{equation}
We now show that the bracketed quantity in~\eqref{eq:total-vel} is exactly
$\nabla(\delta\mathcal{F}/\delta\rho)$.

%==========================================================================
% Step 3: First Variation of the Entropy Term
%==========================================================================
\medskip
\noindent\textbf{Step 3: First variation of the entropy term.}

Consider the relative entropy
\begin{equation}
\Ent(p \mid \rhotrue) = \int p\,\log\frac{p}{\rhotrue}\,dx.
\end{equation}
Take a perturbation $p_{\varepsilon} = p + \varepsilon\eta$ with $\int\eta\,dx = 0$.
A first-order expansion gives
\begin{equation}
\frac{d}{d\varepsilon}\Big|_{\varepsilon=0} \Ent(p_{\varepsilon} \mid \rhotrue)
  = \int \eta(x)\Bigl( \log\frac{p(x)}{\rhotrue(x)} + 1 \Bigr)\,dx.
\end{equation}
Hence the $L^{2}$ first variation is
\begin{equation}
\frac{\delta}{\delta p}\,\Ent(p \mid \rhotrue) = \log\frac{p}{\rhotrue} + 1
\end{equation}
(the additive constant does not affect the gradient flow).
Taking the spatial gradient,
\begin{equation}
\label{eq:ent-grad}
\nabla\frac{\delta}{\delta p}\,\Ent(p \mid \rhotrue)
  = \nabla\log p - \nabla\log\rhotrue.
\end{equation}
Since $\rhotrue \propto e^{-V}$, we have $-\nabla\log\rhotrue = \nabla V$, and therefore
\begin{equation}
\label{eq:ent-var}
\nabla\frac{\delta}{\delta p}\,\Ent(p \mid \rhotrue) = \nabla\log p + \nabla V.
\end{equation}

%==========================================================================
% Step 4: PDE Constraint as a Differentiable Functional
%==========================================================================
\medskip
\noindent\textbf{Step 4: PDE constraint as a differentiable functional.}

The PDE constraint (beyond $\Lburg$) must be expressible as a density
functional~$\mathcal{R}(\rho)$.
Under Hypothesis~3, $\mathcal{R}$ is Fr\'{e}chet-differentiable on $\mathcal{P}_{2}$, so
its first variation $\frac{\delta\mathcal{R}}{\delta\rho}$ exists and satisfies
\begin{equation}
\frac{d}{d\varepsilon}\Big|_{\varepsilon=0} \mathcal{R}(\rho_{\varepsilon})
  = \int \frac{\delta\mathcal{R}}{\delta\rho}(\rho)\,\eta\,dx
\end{equation}
for any zero-mass perturbation $\eta$.

In the Wasserstein setting, perturbations are generated by a transport
field~$\xi$ via the pushforward:
\begin{equation}
T_{\varepsilon}(x) = x + \varepsilon\,\xi(x), \qquad
\rho_{\varepsilon} = (T_{\varepsilon})_{\#}\,\rho.
\end{equation}
The corresponding variational formula is
\begin{equation}
\frac{d}{d\varepsilon}\Big|_{\varepsilon=0}
  \mathcal{R}\bigl((\mathrm{Id} + \varepsilon\xi)_{\#}\rho\bigr)
  = \int \nabla\frac{\delta\mathcal{R}}{\delta\rho}(\rho) \cdot \xi\;\rho\,dx.
\end{equation}
Thus, the $\mathcal{R}$ gradient flow term contributes a drift
$-\nabla(\delta\mathcal{R}/\delta\rho)$ in the velocity field (see~\eqref{eq:guidance-vel}).

%==========================================================================
% Step 5: Identification with the Free Energy Gradient
%==========================================================================
\medskip
\noindent\textbf{Step 5: Identification with the free energy gradient.}

Combining the entropy and PDE-constraint terms, the free energy is
\begin{equation}
\mathcal{F}(\rho) = \Ent(\rho \mid \rhotrue) + \lambda\,\mathcal{R}(\rho).
\end{equation}
From~\eqref{eq:ent-var} and Step~4, its spatial variational gradient is
\begin{equation}
\label{eq:F-grad}
\nabla\frac{\delta\mathcal{F}}{\delta\rho}
  = \nabla\frac{\delta}{\delta\rho}\,\Ent(\rho \mid \rhotrue)
    + \lambda\,\nabla\frac{\delta\mathcal{R}}{\delta\rho}
  = \nabla\log\rho + \nabla V + \lambda\,\nabla\frac{\delta\mathcal{R}}{\delta\rho}.
\end{equation}

Comparing~\eqref{eq:total-vel} with~\eqref{eq:F-grad}, the total effective velocity is exactly
\begin{equation}
\label{eq:vel-F}
v_{\tau}(x) = -a(\tau)\,\nabla_x\frac{\delta\mathcal{F}}{\delta\rho}(\ptau)(x).
\end{equation}

Substituting into the continuity equation (Step~1),
\begin{equation}
\dtau \ptau + \nabla \cdot \Bigl( \ptau\Bigl[ -a(\tau)\,\nabla\frac{\delta\mathcal{F}}{\delta\rho}(\ptau) \Bigr] \Bigr) = 0,
\end{equation}
i.e.,
\begin{equation}
\label{eq:main-pde}
\dtau \ptau = \nabla \cdot \Bigl( a(\tau)\,\ptau\,\nabla\frac{\delta\mathcal{F}}{\delta\rho}(\ptau) \Bigr).
\end{equation}

To eliminate the time-weight $a(\tau)$, define the reparameterized time
\begin{equation}
\theta(\tau) = \int_{0}^{\tau} a(r)\,dr.
\end{equation}
Since $a(\tau) > 0$, $\theta$ is strictly increasing and locally invertible.
By the chain rule,
\begin{equation}
\dtau \ptau = \frac{d\theta}{d\tau}\,\partial_{\theta}p_{\theta}
            = a(\tau)\,\partial_{\theta}p_{\theta}.
\end{equation}
Cancelling $a(\tau)$ yields the canonical Wasserstein-2 gradient flow form
\begin{equation}
\label{eq:W2-final}
\partial_{\theta} p_{\theta}
  = \nabla \cdot \Bigl( p_{\theta}\,\nabla\frac{\delta\mathcal{F}}{\delta p}(p_{\theta}) \Bigr).
\end{equation}

This completes the bridge from the EntroDiff reverse sampler to an exact Wasserstein
gradient flow. $\square$

%==========================================================================
% Step 6: JKO Discretization from the Wasserstein Gradient Flow
%==========================================================================
\medskip
\noindent\textbf{Step 6: JKO discretization from the Wasserstein gradient flow.}

We now show that the implicit Euler discretization of the above gradient flow is equivalent to
the JKO variational scheme.

\smallskip
\noindent\textit{6.1  Implicit Euler discretization.}
Consider a uniform time grid with step $\Delta\theta$ (in reparameterized time).
Given $p^{n} \approx p_{\theta_{n}}$, the implicit Euler discretization of the gradient flow is
\begin{equation}
\label{eq:implicit-euler}
\frac{p^{n+1} - p^{n}}{\Delta\theta}
  = \nabla \cdot \Bigl( p^{n+1}\,\nabla\frac{\delta\mathcal{F}}{\delta\rho}(p^{n+1}) \Bigr).
\end{equation}

\smallskip
\noindent\textit{6.2  Proximal step analogy.}
In Euclidean space, the implicit Euler step for $\dot x = -\nabla F(x)$ is equivalent to the
proximal minimization
\begin{equation}
x^{n+1} = \arg\min_{x} \Bigl\{ F(x) + \frac{1}{2\Delta t}\,\|x - x^{n}\|^{2} \Bigr\}.
\end{equation}
The Wasserstein analogue replaces the Euclidean distance with $\Wass{2}$ and the Euclidean
gradient with the continuity equation structure.
The claim is that
\begin{equation}
\label{eq:jko-claim}
p^{n+1}
  = \arg\min_{p \in \mathcal{P}_{2}(\R^{N})}
    \Bigl\{ \mathcal{F}(p) + \frac{1}{2\Delta\theta}\,\Wass{2}^{2}(p, p^{n}) \Bigr\}
\end{equation}
is the JKO step corresponding to the implicit Euler discretization above.

\smallskip
\noindent\textit{6.3  Variational derivation.}
Define the augmented functional
\begin{equation}
\mathcal{J}(p) = \mathcal{F}(p) + \frac{1}{2\Delta\theta}\,\Wass{2}^{2}(p, p^{n}),
\end{equation}
and let $p^{n+1}$ be a minimizer.
In Wasserstein space, admissible variations are generated by transport maps rather than
additive perturbations.
Define the perturbation
\begin{equation}
T_{\varepsilon}(x) = x + \varepsilon\,\xi(x), \qquad
p_{\varepsilon} = (T_{\varepsilon})_{\#}\,p^{n+1},
\end{equation}
where $\xi$ is an arbitrary smooth, compactly supported vector field.
Optimality requires
\begin{equation}
\frac{d}{d\varepsilon}\Big|_{\varepsilon=0} \mathcal{J}(p_{\varepsilon}) = 0.
\end{equation}

\textbf{Variation of $\mathcal{F}$.}
By the Wasserstein variational formula,
\begin{equation}
\frac{d}{d\varepsilon}\Big|_{\varepsilon=0} \mathcal{F}(p_{\varepsilon})
  = \int \nabla\frac{\delta\mathcal{F}}{\delta\rho}(p^{n+1}) \cdot \xi\;
    p^{n+1}\,dx.
\end{equation}

\textbf{Variation of the $\Wass{2}^{2}$ term.}
Let $\pi$ be the optimal transport plan from $p^{n+1}$ to $p^{n}$, with Kantorovich
potential~$\phi$.
A classical result of optimal transport theory gives
\begin{equation}
\frac{\delta}{\delta\rho}\Bigl( \frac{1}{2}\,\Wass{2}^{2}(\rho, p^{n}) \Bigr) = \phi,
\end{equation}
and therefore
\begin{equation}
\frac{d}{d\varepsilon}\Big|_{\varepsilon=0} \frac{1}{2}\,\Wass{2}^{2}(p_{\varepsilon}, p^{n})
  = \int \nabla\phi \cdot \xi\;
    p^{n+1}\,dx.
\end{equation}

\textbf{First-order condition.}
Combining both terms,
\begin{equation}
\int \Bigl[ \nabla\frac{\delta\mathcal{F}}{\delta\rho}(p^{n+1})
          + \frac{1}{\Delta\theta}\,\nabla\phi \Bigr] \cdot \xi\;
      p^{n+1}\,dx = 0.
\end{equation}
Since $\xi$ is arbitrary, the integrand must vanish $p^{n+1}$-almost everywhere, yielding
\begin{equation}
\label{eq:nabla-phi}
\nabla\phi = -\Delta\theta\,\nabla\frac{\delta\mathcal{F}}{\delta\rho}(p^{n+1}).
\end{equation}

\textbf{Connection to the optimal transport map.}
The Brenier theorem gives the optimal transport map $T$ from $p^{n+1}$ to $p^{n}$ as
$T(x) = x - \nabla\phi(x)$.
Substituting the expression for $\nabla\phi$,
\begin{equation}
T(x) = x + \Delta\theta\,\nabla\frac{\delta\mathcal{F}}{\delta\rho}(p^{n+1}).
\end{equation}

\textbf{Deriving the discrete continuity equation.}
By the pushforward relation $p^{n} = T_{\#}p^{n+1}$, we expand for small $\Delta\theta$
(using the Jacobian expansion of the pushforward):
\begin{equation}
p^{n}(x)
  = p^{n+1}(x)
    - \Delta\theta\,\nabla\cdot\Bigl( p^{n+1}\,\nabla\frac{\delta\mathcal{F}}{\delta\rho}(p^{n+1}) \Bigr)
    + o(\Delta\theta).
\end{equation}
Rearranging,
\begin{equation}
\frac{p^{n+1} - p^{n}}{\Delta\theta}
  = \nabla\cdot\Bigl( p^{n+1}\,\nabla\frac{\delta\mathcal{F}}{\delta\rho}(p^{n+1}) \Bigr) + o(1).
\end{equation}
Taking $\Delta\theta \to 0$, we recover exactly the implicit Euler discretization~\eqref{eq:implicit-euler}
of the Wasserstein gradient flow.

\smallskip
\noindent\textit{6.4  Conclusion of the JKO derivation.}
Thus, a minimizer $p^{n+1}$ of the JKO functional satisfies the implicit Euler discretization
of the Wasserstein gradient flow.
The reverse sampler, whose density evolution was established in Steps~1--5 as the gradient
flow of~$\mathcal{F}$, therefore admits the JKO variational characterization in the
discrete-time limit $\Delta\tau \to 0$.
This completes the proof of Theorem~5.
\end{proof}

%==========================================================================
% Compact Fokker--Planck Expansion
%==========================================================================
\medskip
\noindent\textbf{Compact Fokker--Planck expansion.}

For completeness, expanding the gradient flow structurally: using
$\nabla\log p_{\theta} = \nabla p_{\theta}/p_{\theta}$ and the explicit form of the
entropy variation from Step~3, the gradient flow PDE becomes
\begin{equation}
\begin{aligned}
\partial_{\theta} p_{\theta}
  &= \Lap p_{\theta}
     + \nabla\cdot(p_{\theta}\nabla V)
     + \lambda\,\nabla\cdot\Bigl( p_{\theta}\,\nabla\frac{\delta\mathcal{R}}{\delta\rho} \Bigr),
\end{aligned}
\end{equation}
exhibiting the standard decomposition into a diffusion term ($\Lap p_{\theta}$),
an external potential drift ($\nabla\cdot(p_{\theta}\nabla V)$),
and a PDE-constraint transport term.

%==========================================================================
% Remark: Relation to the Obukhov--Fokker--Planck Correspondence
%==========================================================================
\medskip
\noindent\textbf{Remark (Relation to the Obukhov--Fokker--Planck correspondence).}

The identification~\eqref{eq:vel-F} can be understood through the lens of
entropy-dissipation balances.
The pure reverse probability-flow ODE ($v = -a\nabla\log p$) is the Wasserstein gradient flow
of the negative entropy $H(p) = \int p\log p$.
Adding the PDE guidance terms $\nabla V + \lambda\nabla(\delta\mathcal{R}/\delta\rho)$ shifts
the stationary point from the standard Gaussian to the target~$\rhotrue$ while incorporating
the physical constraint~$\mathcal{R}$.
The resulting free energy $\mathcal{F} = \Ent(\cdot\mid\rhotrue) + \lambda\mathcal{R}$ is
the natural Lyapunov functional for the guided dynamics, and the JKO scheme is its canonical
time discretization~\citep{jko1998}.


%==============================================================================
% Appendix A2 · Extra Experiments
%
% 留作 rebuttal 阶段补充：E4 (shallow-water) / E5 (Vlasov–Poisson)
%==============================================================================
\section{Additional Experiments}
\label{appx:extra-exp}

\subsection{E4 · 2D Shallow-water (rebuttal)}
\label{appx:e4}
\todo{If time permits before rebuttal: hydraulic jump benchmark.}

\subsection{E5 · 1D Vlasov–Poisson (rebuttal)}
\label{appx:e5}
\todo{Cold-beam instability with weak-solution filamentation. Compare against Deep Kinetic JKO \citep{deepkineticjko}.}

\subsection{Extended ablations}
\label{appx:abl}
\todo{Ablate noise schedule, parameterization, loss-family components, and $\lambda$ hyperparameters individually.}

%==============================================================================
% Appendix A3 · Implementation Details
%==============================================================================
\section{Implementation Details}
\label{appx:impl}

\subsection{Network architecture}
\label{appx:impl:arch}
\todo{Detail $\phi^{\mathrm{sm}}$, $\phi^{\mathrm{sh}}$, $\kappa$ networks. Sizes, activations, normalizations.}

\subsection{Training hyperparameters}
\label{appx:impl:train}
\todo{Optimizer (AdamW), batch sizes, learning rates, EMA, training horizon, $\lambda_{\mathrm{ent}}$, $\lambda_{\mathrm{BV}}$, $\lambda_{\mathrm{Burg}}$.}

\subsection{Sampling hyperparameters}
\label{appx:impl:sample}
\todo{Number of Heun steps, $\zeta_{\mathrm{obs}}$, $\zeta_{\mathrm{PDE}}$ schedules.}

\subsection{Compute and reproducibility}
\label{appx:impl:compute}
\todo{GPU type, runtime per benchmark, anonymized code repository link in the supplementary material.}

%==============================================================================

\newpage
\input{checklist.tex}

\end{document}
